% Editorial-Manager pdflatex-compatible supplement wrapper. See main_em.tex
% for the rationale for avoiding fontspec.

\documentclass[preprint,11pt,authoryear,a4paper]{elsarticle}

\usepackage[utf8]{inputenc}
\usepackage[T1]{fontenc}
\usepackage{lmodern}

\usepackage{amsmath,amssymb,amsthm}
\usepackage{booktabs}
\usepackage{graphicx}
\usepackage{siunitx}
\usepackage{xurl}
\usepackage{hyperref}
\hypersetup{hidelinks}
\usepackage[capitalise,nameinlink]{cleveref}
\usepackage{xcolor}
\usepackage{longtable}
\usepackage{enumitem}
\usepackage{array}
\usepackage{calc}
\usepackage{multirow}

\usepackage[a4paper,margin=2.5cm]{geometry}
\usepackage{microtype}
\usepackage[font=small,labelfont=bf,labelsep=period]{caption}
\setlength{\parindent}{1.2em}

% Use \DeclareUnicodeCharacter (inputenc-native) so we don't emit
% "Redefining Unicode character" warnings for §, –, †, ‡ that inputenc
% already handles.
\DeclareUnicodeCharacter{25AB}{\ensuremath{\,\square\,}}
\DeclareUnicodeCharacter{03C1}{\ensuremath{\rho}}
\DeclareUnicodeCharacter{0394}{\ensuremath{\Delta}}
\DeclareUnicodeCharacter{00A7}{\S}
\DeclareUnicodeCharacter{2013}{--}
\DeclareUnicodeCharacter{2020}{\textsuperscript{\dag}}
\DeclareUnicodeCharacter{2021}{\textsuperscript{\ddag}}

\providecommand{\tightlist}{\setlength{\itemsep}{0pt}\setlength{\parskip}{0pt}}
\providecommand{\pandocbounded}[1]{#1}
\providecommand{\textsubscript}[1]{\ensuremath{_{\text{#1}}}}
\providecommand{\real}[1]{#1}

% Corrected 2026-08-15: embed the PNG directly instead of rendering a
% TIFF placeholder box (see manuscript.tex for rationale).
\newcommand{\emfig}[1]{%
  \includegraphics[width=0.92\linewidth]{#1.png}%
}

\renewcommand{\thefigure}{S\arabic{figure}}
\renewcommand{\thetable}{S\arabic{table}}
\renewcommand{\theequation}{S\arabic{equation}}
\renewcommand{\thesection}{S\arabic{section}}

\journal{BioSystems}

\begin{document}

\begin{frontmatter}
\title{Supplementary material: Robust error-minimization in the genetic code in $\mathrm{GF}(2)^6$}
\author[lc]{Paul Clayworth}
\author[bs]{Sergey Kornilov}
\affiliation[lc]{organization={Logocentricity Inc.}, addressline={Austin, TX}, country={USA}}
\affiliation[bs]{organization={Biostochastics, LLC}, addressline={Seattle, WA}, country={USA}}
\end{frontmatter}

\section{Claim hierarchy with full justifications}\label{sec:s-claims}

The 15 evaluated claims are organised into five evidentiary tiers
(Supported / Exploratory / Falsified / Rejected / Tautological) and
listed with full justifications in \Cref{tbl:s-claims}. The hierarchy
is registered in code
(\texttt{src/codon\_topo/reports/claim\_hierarchy.py} in
\url{https://github.com/biostochastics/codontopo}) so that every
analysis run produces a verifiable status report
(\texttt{codon-topo\ claims}) rather than the status table being
maintained as prose. The ``Falsified'' and ``Rejected'' rows record
claims this paper or earlier work has tested and found to be wrong; we
list them so that the framework's negative scope is visible alongside
its positive results.

\begin{longtable}[]{@{}
  >{\raggedright\arraybackslash}p{(\linewidth - 4\tabcolsep) * \real{0.50}}
  >{\centering\arraybackslash}p{(\linewidth - 4\tabcolsep) * \real{0.10}}
  >{\raggedright\arraybackslash}p{(\linewidth - 4\tabcolsep) * \real{0.40}}@{}}
\caption{Complete claim hierarchy with full claim statements and
justifications.}\tabularnewline
\toprule\noalign{}
\begin{minipage}[b]{\linewidth}\raggedright
\textbf{Claim ID and statement}
\end{minipage} & \begin{minipage}[b]{\linewidth}\centering
\textbf{Status}
\end{minipage} & \begin{minipage}[b]{\linewidth}\raggedright
\textbf{Justification}
\end{minipage} \\
\midrule\noalign{}
\endfirsthead
\toprule\noalign{}
\begin{minipage}[b]{\linewidth}\raggedright
\textbf{Claim ID and statement}
\end{minipage} & \begin{minipage}[b]{\linewidth}\centering
\textbf{Status}
\end{minipage} & \begin{minipage}[b]{\linewidth}\raggedright
\textbf{Justification}
\end{minipage} \\
\midrule\noalign{}
\endhead
\bottomrule\noalign{}
\endlastfoot
\begin{minipage}[t]{\linewidth}\raggedright
\textbf{hypercube\_coloring\_optimality}\\
The standard genetic code is significantly error-minimizing under four
established, code-independent physicochemical distance measures with
partially overlapping content: Grantham (\(p = 0.006\)), Miyata
(\(p < 0.001\)), Woese polar requirement (\(p = 0.003\)), and
Kyte--Doolittle hydropathy (\(p = 0.001\)).\strut
\end{minipage} & Supported & Cross-metric sensitivity: Grantham
\(p = 0.006\), Miyata \(p < 0.001\), polar requirement \(p = 0.003\),
Kyte--Doolittle \(p = 0.001\) (quartet-pattern shuffle null,
\(n =\)10,000). Stop penalty sensitivity (0/150/215/300): immaterial. \\
\begin{minipage}[t]{\linewidth}\raggedright
\textbf{per\_table\_optimality\_preservation}\\
26 of 27 NCBI translation tables remain in the top 5\% of their own
quartet-pattern shuffle null for Grantham edge-mismatch (BH--FDR
corrected); only translation table 3 (yeast mito) exceeds the
threshold.\strut
\end{minipage} & Supported & Per-table quartet-pattern shuffle null
applied to all 27 NCBI tables; significant fraction (BH--FDR
\(p < 0.05\)) and per-table quantiles are reported in the main text and
the per-table CSV \texttt{output/tables/T4\_per\_table\_optimality.csv}
(released in the \texttt{codontopo} repository). Translation table 3
(yeast mito) is the marginal exception. \\
\begin{minipage}[t]{\linewidth}\raggedright
\textbf{optimality\_rho\_robustness}\\
Coloring optimality is robust across all diagonal-edge weights
\(\rho \in \lbrack 0,1\rbrack\), including the full Hamming graph
\(H(3,4) = K_{4}\square K_{4}\square K_{4}\) of single-nucleotide
substitutions (\(p < 0.01\) at all \(\rho\) values).\strut
\end{minipage} & Supported & \(\rho\)-sweep at \(n =\)10,000:
\(p \leq 0.006\) at all \(\rho\) values; effect-size \(z\) increases
monotonically from \(\rho = 0\) to \(\rho = 1\). \\
\begin{minipage}[t]{\linewidth}\raggedright
\textbf{topology\_avoidance\_depletion}\\
Natural codon reassignments are depleted for topology-breaking changes
(moves that fragment an amino acid's codon family) at approximately 21\%
of observed events vs 66--73\% of the candidate landscape, robust to
adjacency definition (\(Q_{6}\) vs the full Hamming graph \(H(3,4)\))
and clade exclusion.\strut
\end{minipage} & Supported & Permutation \(p \leq 10^{- 4}\) under both
\(Q_{6}\) and \(H(3,4)\); hypergeometric \(p = 1.6 \times 10^{- 8}\)
(\(Q_{6}\), new disconnection) and \(p = 1.3 \times 10^{- 6}\)
(\(H(3,4)\), \(\Delta\beta_{0} > 0\)). 5--7 of 28 de-duplicated events
are topology-breaking versus \(\approx 64\)--\(75\)\% of the candidate
landscape (\(\approx 3.0\)--\(3.4\)-fold depletion). \(Q_{6}\) is
encoding-dependent (8 of 24 bijections give no depletion); \(H(3,4)\) is
encoding-independent and is the primary test. Clade-exclusion
sensitivity (7 regimes): all \(p < 10^{- 5}\). \\
\begin{minipage}[t]{\linewidth}\raggedright
\textbf{trna\_enrichment\_reassigned\_aa}\\
Organisms with variant genetic codes tend to show elevated tRNA gene
copy numbers for the reassigned amino acid, but the enrichment does not
attain a strict worst-case-subset threshold and the
topology-breaking-only subset is null.\strut
\end{minipage} & Exploratory & 332 maximal independent sets of size 6
(Bron--Kerbosch on the complement of the conflict graph): median
Stouffer \(p = 0.037\), best \(p = 0.012\), worst \(p = 0.104\); 264/332
(79.5\%) sets fall below 0.05. Topology-breaking subset (\(n = 4\))
Stouffer \(p = 0.387\) (null). 18 tRNAscan-SE-verified assemblies (15
variant + 3 standard controls), 24 pairings across 5 variant codes. \\
\begin{minipage}[t]{\linewidth}\raggedright
\textbf{bit\_position\_bias\_weighted}\\
Codon reassignment bit-flip distribution shows positional skew in
\({\text{GF}(2)}^{6}\) coordinates under a uniform null, but the signal
does not survive nulls that respect the source-codon non-independence of
recurrent reassignments. Reported as exploratory/null rather than a
positive finding.\strut
\end{minipage} & Exploratory & Uniform null \(p = 0.006\) but inflated
by recurrent-reassignment non-independence (e.g., the recurrent
UGA\(\rightarrow\)Trp event contributes the same source codon across
many lineages). De-duplicating to unique (codon, target) pairs gives
\(p = 0.075\). The codon-preserving null (which permutes target amino
acids while holding the source codon fixed and is the appropriate
non-independence-respecting null for this question) absorbs the apparent
bias entirely. We retain the row in the registry as a cautionary record
of how the choice of null can flip a putative positive into a null
finding, and we treat the bit-position-bias claim as not supported. \\
\begin{minipage}[t]{\linewidth}\raggedright
\textbf{mechanism\_boundary\_conditions}\\
tRNA gene duplication accompanies codon reassignment in large nuclear
genomes (ciliates, yeasts) but not in streamlined genomes
(\emph{Blastocrithidia}: anticodon stem shortening;
\emph{Mycoplasmoides}: anticodon modification).\strut
\end{minipage} & Exploratory & Three-tier pattern: duplication / stem
shortening / modification. Descriptive. \\
\begin{minipage}[t]{\linewidth}\raggedright
\textbf{atchley\_f3\_serine\_convergence}\\
Serine's extreme Atchley Factor 3 score (\(F_{3} = - 4.760\), most
extreme of 20 amino acids, 2.24 SD below mean) converges with the
\({\text{GF}(2)}^{6}\) topological disconnection: \(F_{3}\) captures the
mismatch between Serine's small physicochemical footprint and its
disproportionate codon diversity, which the geometric framework
identifies as maximal inter-family Hamming distance among 6-codon amino
acids (4 vs 1 for Leu and Arg).\strut
\end{minipage} & Exploratory & Serine \(F_{3} = - 4.760\), 2.24 SD below
mean. Complementary, not independent. \\
\begin{minipage}[t]{\linewidth}\raggedright
\textbf{variant\_code\_disconnection\_catalogue}\\
Systematic survey of Hamming-graph connectivity across 27 NCBI
translation tables identifies 4 lineage-collapsed variant-code
amino-acid disconnections at \(\varepsilon = 1\) under the default
encoding: Thr in yeast mitochondrial code (table 3), Leu in
chlorophycean mitochondrial codes (tables 16 and 22, collapsed to a
single algal-mito event), Ala in \emph{Pachysolen tannophilus} nuclear
code (table 26), and tripartite Ser in \emph{Candida}-clade alternative
yeast nuclear code (table 12).\strut
\end{minipage} & Exploratory & 4 variant-code cases at
\(\varepsilon = 1\) in \(Q_{6}\) (default encoding): Thr (Table 3, yeast
mito), Leu (Tables 16/22, chlorophycean and \emph{S. obliquus} mito;
collapsed to one algal-mito event), Ala (Table 26, \emph{Pachysolen}),
Ser (Table 12, \emph{Candida}). Separately, Table 32 (Balanophoraceae
plastid, UAG\ensuremath{\to}{}Trp) creates a 2-fold Trp pair (UGG, UAG) that remains
connected at \(\varepsilon\)=1 but breaks the bit-5 two-fold filtration (a filtration
finding, not a disconnection). \\
\begin{minipage}[t]{\linewidth}\raggedright
\textbf{kras\_fano\_clinical\_prediction}\\
XOR ``Fano'' relationships in \({\text{GF}(2)}^{6}\) predict enrichment
of specific amino acids at KRAS G12 co-mutation sites.\strut
\end{minipage} & Falsified & \(p = 1.0\) across all 6 G12 variants.
\(n =\)1,670 MSK-IMPACT mutations. \\
\begin{minipage}[t]{\linewidth}\raggedright
\textbf{serine\_min\_distance\_4\_invariant}\\
Serine's minimum UCN--AGY Hamming distance \(= 4\) holds invariantly
across all 24 base-to-bit encodings.\strut
\end{minipage} & Rejected & 16/24 encodings give distance 2. Only 8/24
give distance 4. \\
\begin{minipage}[t]{\linewidth}\raggedright
\textbf{psl\_2\_7\_symmetry}\\
\(\text{PSL}(2,7)\) is the fundamental symmetry group of the genetic
code.\strut
\end{minipage} & Rejected & No 64-dim irrep. \citep{antoneli2011}. \\
\begin{minipage}[t]{\linewidth}\raggedright
\textbf{holomorphic\_embedding}\\
The coordinate-wise map
\({\text{GF}(2)}^{6} \rightarrow {\mathbb{C}}^{3}\) sending base-pairs
to fourth roots of unity is a holomorphic embedding extending a
character of \({\text{GF}(8)}^{\ast}\).\strut
\end{minipage} & Rejected & Domain is finite discrete. Character
identity fails: \(i^{2} = - 1 \neq 1\). \\
\begin{minipage}[t]{\linewidth}\raggedright
\textbf{two\_fold\_bit\_5\_filtration}\\
All 2-fold synonymous codon pairs in the standard code differ at exactly
bit position 5 under the default encoding.\strut
\end{minipage} & Tautological & Forced by encoding choice. Holds in
16/24 encodings. \\
\begin{minipage}[t]{\linewidth}\raggedright
\textbf{four\_fold\_prefix\_filtration}\\
All 4-fold synonymous codon groups share a 4-bit prefix with 2-bit
suffixes exhausting \({\text{GF}(2)}^{2}\).\strut
\end{minipage} & Tautological & Trivial under any bijection from 4 bases
to \({\text{GF}(2)}^{2}\). \\
\end{longtable}

\protect\phantomsection\label{tbl:s-claims}{}

\section{Encoding sensitivity analysis (coloring
optimality)}\label{sec:s-encoding}

There are \(4! = 24\) distinct bijections from
\(\left\{ C,U,A,G \right\}\) to \({\text{GF}(2)}^{2}\) (24 ways to
assign the four nucleotide letters to the four 2-bit binary patterns
00/01/10/11). The default encoding used in the main text is
\(C \mapsto 00\), \(U \mapsto 01\), \(A \mapsto 10\), \(G \mapsto 11\)
(rationale stated in main-text §2.1). To check that the
coloring-optimality result is not specific to that choice, we re-ran the
Grantham quartet-pattern shuffle null analysis under each of the 23
alternative encodings, holding everything else fixed (per-encoding
\(n =\)10,000 null draws, seed 135325). All \{'n\_encodings': 24,
`null\_type': `freeland\_hurst', `quantile\_min': 0.01, `quantile\_max':
2.77, `mean\_quantile': 1.0929166666666668, `all\_significant\_p05':
True, `p\_max': 0.027797220277972202, `p\_min':
0.00019998000199980003\}.at(``n\_encodings'', default: 24) encodings
yield significant optimality (\(p < 0.05\) under the quartet-pattern
shuffle null), with all per-encoding conservative \(p\)-values bracketed
by \(0.0002\) (best) and \(0.0278\) (worst). The mean per-encoding
quantile is 1.09\% (range 0.01\%--2.77\%), confirming that the result is
not an artifact of the default encoding.

Properties that are encoding-invariant:

\begin{itemize}
\item
  Serine disconnection at \(\varepsilon = 1\) (holds under all 24
  encodings)
\item
  Coloring optimality significance (all 24 significant)
\item
  Four-fold prefix filtration (tautological under any bijection)
\item
  \(H(3,4)\) adjacency graph (encoding-independent by construction)
\end{itemize}

Properties that are encoding-dependent:

\begin{itemize}
\item
  Serine inter-family minimum Hamming distance (4 in 8/24 encodings, 2
  in 16/24)
\item
  Two-fold bit-5 filtration (holds in 16/24 encodings)
\item
  \(Q_{6}\) (Hamming-1) adjacency partition of \(H(3,4)\), and the
  topology-avoidance signal computed under \(Q_{6}\) adjacency (see §S4)
\item
  Specific score values and rank orderings
\end{itemize}

Full per-encoding results are written to
\texttt{output/coloring\_optimality.json} (block
\texttt{encoding\_sensitivity}) by the standard \texttt{codon-topo\ all}
invocation, and reproducible from the \texttt{codontopo} repository.

\begin{figure}
\centering
\emfig{FigA\_coloring\_null}
\caption{Hypercube coloring null distribution (extended). Histogram of
Grantham edge-mismatch scores \(F\) across the quartet-pattern shuffle
null draws under the default encoding (C=00, U=01, A=10, G=11); the
observed standard-code score is marked and its quantile shown. Companion
to main-text Figure 2A, retained here at higher resolution together with
the per-encoding sweep and per-table breakdown below.}
\end{figure}

\protect\phantomsection\label{fig:s-coloring-null}{}

\begin{figure}
\centering
\emfig{FigD\_score\_decomposition}
\caption{Grantham mismatch score decomposition by nucleotide position.
Contribution of each codon position (1, 2, 3) to the total Grantham
edge-mismatch score \(F\) under the default encoding. Position 2
dominates, consistent with the biochemical hierarchy of mutational
impact. This is a descriptive breakdown of the summed score, not an
inferential test.}
\end{figure}

\protect\phantomsection\label{fig:s-score-decomposition}{}

\subsection{Null-model choice: quartet-pattern vs classical Haig-Hurst
AA-permutation}\label{sec:s-null-hh-aa}

The primary coloring-optimality tests throughout this paper use the
\textbf{quartet-pattern shuffle} null (main-text §2.3.1: the 64 codons
are grouped into 16 first-two-base quartets, and the amino-acid pattern
at each quartet is held as an atomic labeled 4-tuple; the two
stop-containing quartets in the standard code (UA and UG) are held fixed
and the remaining 14 patterns are permuted uniformly at random among the
14 non-stop quartet slots). This is distinct from the classical
Haig--Hurst 1991 / Freeland--Hurst 1998 null commonly cited in the
code-optimization literature, which fixes the standard code's 20
sense-codon families (the unlabeled partition of the 61 sense codons
into synonymous groups) and permutes the 20 amino-acid labels uniformly
across those families with stop positions held fixed. The two ensembles
preserve different structural properties
(\Cref{tbl:s-null-preservation}) and admit different numbers of
alternative codes (\(\approx 14!\) for the quartet-pattern shuffle
versus \(\approx 20!\) for the AA-permutation null). The
quartet-pattern-shuffle count is \(14!\) rather than \(16!\) because in
the standard code both stop-containing quartets (UA and UG) are held
fixed, leaving \(14\) mobile quartet slots.

\begin{longtable}[]{@{}
  >{\raggedright\arraybackslash}p{(\linewidth - 4\tabcolsep) * \real{0.4444}}
  >{\centering\arraybackslash}p{(\linewidth - 4\tabcolsep) * \real{0.2778}}
  >{\centering\arraybackslash}p{(\linewidth - 4\tabcolsep) * \real{0.2778}}@{}}
\caption{\textbf{Preservation properties of the two null ensembles.} The
two nulls preserve non-nested structural properties. The quartet-pattern
shuffle preserves each named amino acid's total codon count (because the
atomic 4-tuples that move carry their internal AA composition with them)
but permutes which specific codons decode to a given amino acid; the
classical Haig--Hurst AA-permutation preserves which codons decode
together (the family partition) but allows a named amino acid's codon
count to change (Trp's single-codon family may become a 6-codon family
if its label is swapped with Leu's). Neither ensemble is uniformly
``more stringent.'' The quartet-pattern-shuffle state space is \(14!\)
rather than \(16!\) because in the standard code both stop-containing
quartets (UA, containing UAA/UAG, and UG, containing UGA) are held
fixed, leaving 14 mobile quartet slots.}\tabularnewline
\toprule\noalign{}
\begin{minipage}[b]{\linewidth}\raggedright
\textbf{Property}
\end{minipage} & \begin{minipage}[b]{\linewidth}\centering
\textbf{Quartet-pattern shuffle}
\end{minipage} & \begin{minipage}[b]{\linewidth}\centering
\textbf{Classical Haig--Hurst AA-permutation}
\end{minipage} \\
\midrule\noalign{}
\endfirsthead
\toprule\noalign{}
\begin{minipage}[b]{\linewidth}\raggedright
\textbf{Property}
\end{minipage} & \begin{minipage}[b]{\linewidth}\centering
\textbf{Quartet-pattern shuffle}
\end{minipage} & \begin{minipage}[b]{\linewidth}\centering
\textbf{Classical Haig--Hurst AA-permutation}
\end{minipage} \\
\midrule\noalign{}
\endhead
\bottomrule\noalign{}
\endlastfoot
Stop-codon positions & fixed & fixed \\
Codon-family partition (which codons decode to a common amino acid, as
an unlabeled clustering) & permuted & fixed \\
Per-named-amino-acid codon count (e.g., Leu's total codon count) & fixed
& permuted \\
Sorted degeneracy multiset (\((6,6,6,4,4,\ldots,1)\)) & fixed & fixed \\
Quartet-slot AA-pattern shape at each first-two-base slot & permuted
(atomic 4-tuples redistribute across slots) & fixed (labels within each
slot swap only) \\
Approx. number of possible codes & \(14! \approx 8.7 \times 10^{10}\) &
\(20! \approx 2.4 \times 10^{18}\) \\
Historical reference & not standard; introduced here &
\citep{haig1991, freeland1998} \\
\end{longtable}

\protect\phantomsection\label{tbl:s-null-preservation}{}

We reran the four physicochemical metrics under the classical
Haig--Hurst AA-permutation null (\(n =\)10,000 draws, seed 135325) for
direct comparison against the quartet-pattern shuffle values reported in
main-text Table 2. Results are in \Cref{tbl:s-null-comparison}.

\begin{longtable}[]{@{}lrrrrrr@{}}
\caption{\textbf{Coloring-optimality \(p\)-values under the
quartet-pattern shuffle (QP, primary) vs the classical Haig--Hurst
AA-permutation null (HH-AA, sensitivity).} Both use \(n =\)10,000 draws,
seed 135325. Conservative \(p\)-values are \((k + 1)/(n + 1)\) with the
tail convention \(f \leq F_{\text{obs}}\) (draws as extreme as or more
extreme than the observed score). Treating the eight reported
\(p\)-values as one test family, all eight pass Bonferroni at
\(\alpha = 0.05/8 = 0.00625\). For three of the four code-independent
metrics the classical HH-AA null gives a \textbf{more} extreme \(p\):
because the HH-AA ensemble admits a wider space of alternative codes
(\(20!\) vs \(14!\)), the standard code sits further into its low-cost
tail. Kyte--Doolittle inverts: hydropathy is unusually well-aligned with
the wobble-quartet structure that QP preserves, so the QP null contains
many low-hydropathy-mismatch alternatives that the HH-AA null (which
breaks that alignment) does not, and the observed code's percentile is
deeper against QP than against HH-AA. The conclusion --- that the
standard code is significantly error-minimizing across all four
code-independent physicochemical metrics --- is robust to the choice of
null.}\tabularnewline
\toprule\noalign{}
\textbf{Metric} & \textbf{Observed \(F\)} & \textbf{QP null \(z\)} &
\textbf{QP \(p\)} & \textbf{HH-AA null \(z\)} & \textbf{HH-AA \(p\)} &
\textbf{Δ (HH \ensuremath{-} QP)} \\
\midrule\noalign{}
\endfirsthead
\toprule\noalign{}
\textbf{Metric} & \textbf{Observed \(F\)} & \textbf{QP null \(z\)} &
\textbf{QP \(p\)} & \textbf{HH-AA null \(z\)} & \textbf{HH-AA \(p\)} &
\textbf{Δ (HH \ensuremath{-} QP)} \\
\midrule\noalign{}
\endhead
\bottomrule\noalign{}
\endlastfoot
Grantham & 13477 & 2.35 & 0.0062 & 2.76 & 0.0021 & HH more extreme \\
Miyata & 235.3 & 3.14 & 0.0004 & 4.13 & 0.0001 & HH more extreme \\
Polar requirement & 347.3 & 2.73 & 0.0026 & 3.45 & 0.0001 & HH more
extreme \\
Kyte--Doolittle & 457.2 & 2.78 & 0.0014 & 2.68 & 0.0027 & QP more
extreme \\
\end{longtable}

\protect\phantomsection\label{tbl:s-null-comparison}{}

\textbf{Why we do not describe either null as ``more stringent'' than
the other.} The two ensembles preserve non-nested structural properties
(\Cref{tbl:s-null-preservation}), and empirical tail behaviour
(\Cref{tbl:s-null-comparison}) shows the metric-specific direction of
the \(p\)-value shift is not monotone across the four metrics. We
therefore report both, note the null under which each headline
\(p\)-value was computed (the QP shuffle), and treat the HH-AA null as a
same-direction sensitivity check rather than as a validation or a
refutation of the QP tail.

\section{Topology-breaking definitions: 2 \ensuremath{\times} 2
audit}\label{sec:s-topology-defs}

The topology-avoidance result depends on two independent definitional
choices: (i) which adjacency graph defines codon-family connectivity
(\(Q_{6}\) Hamming-1 in the default \({\text{GF}(2)}^{6}\) encoding, vs
the encoding-independent \(H(3,4)\) single-nucleotide graph), and (ii)
what counts as a ``topology-breaking'' candidate move (a candidate move
creates a \textbf{new} disconnection in a previously connected family,
vs the candidate move strictly increases \(\beta_{0}\) summed across
families). To prevent silent dependence on either choice, we report all
four combinations.

We define two notions of ``topology-breaking'' candidate move:

\begin{enumerate}
\item
  \textbf{New disconnection in a previously connected family}: a
  candidate move makes some amino acid disconnected at
  \(\varepsilon = 1\) that was connected in the standard code (i.e., the
  amino acid was not previously in the disconnection catalogue but is
  after the move).
\item
  \textbf{Increase in components} (\(\Delta\beta_{0} > 0\), the
  conditional-logit feature): the total number of connected components,
  summed across amino-acid codon graphs, strictly increases:
  \[\Delta_{\text{topo }} = \sum_{a}\beta_{0}\left( G_{a}^{\text{after}} \right) - \sum_{a}\beta_{0}\left( G_{a}^{\text{before}} \right) > 0.\]
\end{enumerate}

Both definitions are reported under both \(Q_{6}\) adjacency (Hamming-1
in the default \({\text{GF}(2)}^{6}\) encoding) and \(H(3,4)\) adjacency
(full single-nucleotide adjacency, encoding-independent), giving four
cells. All four share the same denominators (1,280 candidate moves, 28
de-duplicated observed events). The full \(2 \times 2\) result is shown
in \Cref{tbl:s-topo-audit}.

\begin{longtable}[]{@{}
  >{\raggedright\arraybackslash}p{(\linewidth - 10\tabcolsep) * \real{0.1667}}
  >{\raggedright\arraybackslash}p{(\linewidth - 10\tabcolsep) * \real{0.1667}}
  >{\raggedleft\arraybackslash}p{(\linewidth - 10\tabcolsep) * \real{0.1667}}
  >{\raggedleft\arraybackslash}p{(\linewidth - 10\tabcolsep) * \real{0.1667}}
  >{\raggedleft\arraybackslash}p{(\linewidth - 10\tabcolsep) * \real{0.1667}}
  >{\raggedleft\arraybackslash}p{(\linewidth - 10\tabcolsep) * \real{0.1667}}@{}}
\caption{\textbf{Topology-avoidance \(2 \times 2\) definition \(\times\)
adjacency audit.} All four cells share the same denominators (1,280
candidate moves, 28 de-duplicated observed events) and differ only in
(i) which adjacency graph defines codon-family connectivity (\(Q_{6}\)
Hamming-1 vs encoding-independent \(H(3,4)\)) and (ii) what counts as a
topology-breaking move (\(\Delta\beta_{0} > 0\) summed over amino acids
vs creation of a \textbf{new} disconnection in a previously connected
family). \(K/N\) = topology-breaking candidates among all candidates;
\(x/n\) = topology-breaking observed events among all observed events.
All four cells yield depletion in the same direction with comparable
risk ratios (0.28--0.33) and hypergeometric \(p < 10^{- 5}\). The main
text uses the \(H(3,4)\), \(\Delta\beta_{0} > 0\) cell as primary
(encoding-independent adjacency, definition matched to the
conditional-logit feature).}\tabularnewline
\toprule\noalign{}
\begin{minipage}[b]{\linewidth}\raggedright
\textbf{Adjacency}
\end{minipage} & \begin{minipage}[b]{\linewidth}\raggedright
\textbf{Topology-breaking definition}
\end{minipage} & \begin{minipage}[b]{\linewidth}\raggedleft
\textbf{\(K/N\)}
\end{minipage} & \begin{minipage}[b]{\linewidth}\raggedleft
\textbf{\(x/n\)}
\end{minipage} & \begin{minipage}[b]{\linewidth}\raggedleft
\textbf{Hyper. \(p\)}
\end{minipage} & \begin{minipage}[b]{\linewidth}\raggedleft
\textbf{RR (95\% CI)}
\end{minipage} \\
\midrule\noalign{}
\endfirsthead
\toprule\noalign{}
\begin{minipage}[b]{\linewidth}\raggedright
\textbf{Adjacency}
\end{minipage} & \begin{minipage}[b]{\linewidth}\raggedright
\textbf{Topology-breaking definition}
\end{minipage} & \begin{minipage}[b]{\linewidth}\raggedleft
\textbf{\(K/N\)}
\end{minipage} & \begin{minipage}[b]{\linewidth}\raggedleft
\textbf{\(x/n\)}
\end{minipage} & \begin{minipage}[b]{\linewidth}\raggedleft
\textbf{Hyper. \(p\)}
\end{minipage} & \begin{minipage}[b]{\linewidth}\raggedleft
\textbf{RR (95\% CI)}
\end{minipage} \\
\midrule\noalign{}
\endhead
\bottomrule\noalign{}
\endlastfoot
\(H(3,4)\) (primary) & \(\Delta\beta_{0} > 0\) (increase in components)
& 846 / 1,280 & 6 / 28 & \(1.3 \times 10^{- 6}\) & 0.32 (0.16--0.66) \\
\(H(3,4)\) & new disconnection in connected family & 822 / 1,280 & 5 /
28 & \(5.0 \times 10^{- 7}\) & 0.28 (0.13--0.62) \\
\(Q_{6}\) & \(\Delta\beta_{0} > 0\) (increase in components) & 963 /
1,280 & 7 / 28 & \(2.2 \times 10^{- 8}\) & 0.33 (0.17--0.63) \\
\(Q_{6}\) & new disconnection in connected family & 931 / 1,280 & 6 / 28
& \(1.6 \times 10^{- 8}\) & 0.29 (0.14--0.60) \\
\end{longtable}

\protect\phantomsection\label{tbl:s-topo-audit}{}

The full machine-readable audit is released as part of the
\texttt{codontopo} repository's \texttt{output/} artifacts (the
\texttt{definitions\_audit} block in the topology-avoidance results
file).

\begin{figure}
\centering
\emfig{FigF\_topology\_avoidance}
\caption{Topology-avoidance depletion under the primary \(H(3,4)\)
adjacency, \(\Delta\beta_{0} > 0\) definition (row 1 of
\Cref{tbl:s-topo-audit}). Observed rate of topology-breaking events
(bar with permutation \(p\)-value annotation) contrasted with the
candidate-landscape rate (dashed reference line). All four cells of the
\(2 \times 2\) audit above give qualitatively equivalent depletion (RR
0.28--0.33, hypergeometric \(p < 10^{- 5}\)); this figure visualises the
primary cell.}
\end{figure}

\protect\phantomsection\label{fig:s-topology-avoidance}{}

\section{\texorpdfstring{Topology avoidance: \(Q_{6}\) encoding-sweep
sensitivity}{Topology avoidance: Q\_\{6\} encoding-sweep sensitivity}}\label{sec:s-encoding-sweep}

This section documents what motivated the demotion of \(Q_{6}\) from
primary to secondary topology adjacency in the manuscript.

The \(H(3,4)\) Hamming graph is encoding-independent: every two-bit
bijection from \(\left\{ A,C,G,U \right\}\) to
\(\left\{ 0,1 \right\}^{2}\) produces the same nucleotide-level
adjacency, because \(H(3,4)\) depends only on which nucleotide letters
differ between two codons, not on their binary encoding. The \(Q_{6}\)
subgraph, however, depends on the encoding because the partition into
Hamming-1 (single-bit-change) edges versus Hamming-2 (within-nucleotide
diagonal) edges is bijection-specific: under one encoding two codons
that differ at exactly one nucleotide may project to a Hamming-1 edge of
\(Q_{6}\), and under another encoding they may project to a Hamming-2
edge. To test whether the \(Q_{6}\) topology-avoidance result survives
this representation choice, we recomputed the \(Q_{6}\)
candidate-landscape rate, observed rate, depletion fold, and
hypergeometric \(p\)-value under all 24 base-to-bit bijections, holding
the same 1,280 candidate moves and 28 observed events.

Across all 24 encodings:

\begin{itemize}
\item
  candidate-landscape rate: min 36\%, median 69.9\%, max 72.7\%
\item
  observed rate: min 21.4\%, median 28.6\%, max 35.7\%
\item
  depletion fold: min 1.01, median 2.45, max 3.39
\item
  hypergeometric \(p\): min \(1.58 \times 10^{- 8}\), median
  \(6.08 \times 10^{- 6}\), max 0.572
\end{itemize}

The default encoding (C=00, U=01, A=10, G=11) gives the largest
depletion and the smallest \(p\). Eight of 24 encodings give a
candidate-landscape rate of approximately 36\%, under which the observed
rate of 21--36\% does not significantly differ from candidate
(\(p > 0.5\)). The \(Q_{6}\) result is therefore not encoding-invariant.
We accordingly present the encoding-independent \(H(3,4)\) result as the
primary topology-avoidance test, with \(Q_{6}\) reported as a
representation-specific decomposition for continuity with the broader
\({\text{GF}(2)}^{6}\) framework. The default encoding was adopted in
companion methodological work (\citep{clayworth2026}) prior to the
present encoding sweep, on the grounds of visualization clarity (it
places the standard code's nine 2-fold-degenerate amino acids on bit-5
differences). We make no claim that this encoding is biologically
privileged. The full per-encoding sweep is released as part of the
\texttt{codontopo} repository's \texttt{output/} artifacts (the
\texttt{Q6\_encoding\_sweep} block in the topology-avoidance results
file); \Cref{fig:s-encoding-sweep} visualizes the per-encoding
depletion fold and hypergeometric \(p\).

\begin{figure}
\centering
\emfig{FigS\_encoding\_sweep}
\caption{Per-encoding \(Q_{6}\) topology-avoidance depletion across all
24 base-to-bit bijections. Bars are sorted by depletion fold; the dashed
line at \(1.0\) marks no-depletion. Bars highlighted in blue are
statistically significant (\(p < 0.05\)); grey bars are not. Eight of 24
encodings (depletion fold \(\approx 1.0\), \(p > 0.5\)) place the
\(Q_{6}\) candidate-landscape rate at \(\approx 36\)\% rather than 73\%,
eliminating the depletion signal. The default encoding (C=00, U=01,
A=10, G=11) yields the largest depletion (3.4-fold). The
encoding-independent \(H(3,4)\) result is constant across all 24
bijections and is the primary topology-avoidance test in this work.}
\end{figure}

\protect\phantomsection\label{fig:s-encoding-sweep}{}

\begin{figure}
\centering
\emfig{FigC\_rho\_robustness}
\caption{Coloring optimality across the \(\rho\)-interpolation between
\(Q_{6}\) (Hamming-1 only, \(\rho = 0\)) and \(H(3,4)\) (full
single-nucleotide mutation graph, \(\rho = 1\)).
\(F_{\rho} = F_{H_{1}} + \rho \cdot F_{H_{2}}\) is the
diagonal-edge-weighted mismatch score; a lower \(F_{\rho}\)-percentile
indicates stronger optimality. The standard code retains top-percentile
placement across the full \(\rho \in \lbrack 0,1\rbrack\) range under
the default encoding, so the coloring-optimality result is not a
knife-edge property of a specific weighting of the two edge classes.
Companion to main-text §3.2.}
\end{figure}

\protect\phantomsection\label{fig:s-rho-robustness}{}

\section{Reassignment candidate-universe denominator
sensitivity}\label{sec:s-denominator}

The hypergeometric and permutation tests of topology avoidance, and the
conditional-logit denominator, both depend on a definition of the
candidate-move universe \(\mathcal{M}(C)\) from the standard code \(C\).
Two definitional choices interact: whether stop-codon \textbf{targets}
are admitted (\(y \in \mathcal{A}_{20}\) vs
\(y \in \mathcal{A}_{20} \cup \left\{ \text{Stop} \right\}\)), and
whether identity moves (\(y = C(x)\)) are excluded. The four
combinations give the variants below; we adopt U1 as primary throughout
and report U2 and U4 as sensitivity universes. All four universes admit
stop \textbf{sources} (i.e. \(x \in \mathcal{C}\) for all 64 codons);
the earlier framing that referred to a source-stop axis was inaccurate
and has been removed.

Formally, the primary candidate-universe is
\(\mathcal{M}(C) = \left\{ (x,y):x \in \mathcal{C},y \in \mathcal{A}_{20} \cup \left\{ \text{Stop} \right\},y \neq C(x) \right\}\)
with \(\left| {\mathcal{M}(C)} \right| = 64 \times 20 =\)1,280. The four
definitional variants and their sizes are listed in
\Cref{tbl:s-denominators}.

\begin{longtable}[]{@{}
  >{\raggedright\arraybackslash}p{(\linewidth - 4\tabcolsep) * \real{0.3333}}
  >{\raggedleft\arraybackslash}p{(\linewidth - 4\tabcolsep) * \real{0.3333}}
  >{\raggedright\arraybackslash}p{(\linewidth - 4\tabcolsep) * \real{0.3333}}@{}}
\caption{\textbf{Candidate-universe denominator definitions for the
topology-avoidance test.} The four variants (U1--U4) interact along
three definitional axes: whether stop-codon targets are admitted
(\(\mathcal{A}_{20} \cup \left\{ \text{Stop} \right\}\) vs
\(\mathcal{A}_{20}\)), whether identity moves \(y = C(x)\) are excluded,
and whether the source codon may itself be a stop codon. We adopt U1 as
the primary universe because it (i) has a uniform alternative count per
codon, (ii) cleanly excludes identity moves which contribute no signal,
and (iii) admits stop-codon reassignment, which is biologically
attested.}\tabularnewline
\toprule\noalign{}
\begin{minipage}[b]{\linewidth}\raggedright
\textbf{Universe}
\end{minipage} & \begin{minipage}[b]{\linewidth}\raggedleft
\textbf{\(\left| \mathcal{M} \right|\)}
\end{minipage} & \begin{minipage}[b]{\linewidth}\raggedright
\textbf{Definition}
\end{minipage} \\
\midrule\noalign{}
\endfirsthead
\toprule\noalign{}
\begin{minipage}[b]{\linewidth}\raggedright
\textbf{Universe}
\end{minipage} & \begin{minipage}[b]{\linewidth}\raggedleft
\textbf{\(\left| \mathcal{M} \right|\)}
\end{minipage} & \begin{minipage}[b]{\linewidth}\raggedright
\textbf{Definition}
\end{minipage} \\
\midrule\noalign{}
\endhead
\bottomrule\noalign{}
\endlastfoot
U1: 21-label, no identity (primary) & 1,280 &
\(y \in \mathcal{A}_{20} \cup \left\{ \text{Stop} \right\}\),
\(y \neq C(x)\); each codon has exactly 20 alternative labels \\
U2: AA-only, no identity & 1,219 & \(y \in \mathcal{A}_{20}\),
\(y \neq C(x)\); 61 sense codons get 19 AA alternatives, 3 stop codons
get 20 AA alternatives \\
U3: AA, with no-ops & 1,280 & \(y \in \mathcal{A}_{20}\), no identity
restriction; each of 64 codons gets 20 AA candidates. Because
\(y \in \mathcal{A}_{20}\) excludes Stop, the 3 stop codons have no AA
identity to match, so only the 61 sense codons contribute an AA no-op
(\(y = C(x)\)); the row total 1,280 = 61 no-ops + 1,219 changes. \\
U4: stop-inclusive with no-ops & 1,344 &
\(y \in \mathcal{A}_{20} \cup \left\{ \text{Stop} \right\}\), no
identity restriction; 64 no-ops included \\
\end{longtable}

\protect\phantomsection\label{tbl:s-denominators}{}

Topology-avoidance results under U2 and U4 are qualitatively identical
(depletion remains \(p < 10^{- 5}\)); the per-cell counts differ only by
the constant rescaling
\(K_{2} = K_{1} - \text{(stop-target candidates)}\) for U2.
\textbf{Observed sense-to-Stop events under U2:} the observed
reassignment corpus contains several sense-to-Stop events (e.g. NCBI
translation table 22 UCA\(\rightarrow\)Stop in \emph{Scenedesmus
obliquus} mitochondrial). Because U2 excludes Stop targets, these events
are \textbf{not eligible} under U2 and are dropped from both the
observed numerator and the candidate denominator of the U2
hypergeometric; the U2 test therefore evaluates whether the
sense-to-sense subset of natural reassignments is topology-depleted
against the sense-to-sense candidate landscape. Sense-to-Stop events
remain in U1 and U4. Detailed per-universe numbers, including the exact
observed counts kept or dropped under each variant, are released as part
of the \texttt{codontopo} repository's \texttt{output/} artifacts (the
\texttt{denominator\_sensitivity} block in the topology-avoidance
results file).

\section{Conditional logit: IIA assumption and explanatory
framing}\label{sec:s-iia}

The conditional-logit framework assumes Independence of Irrelevant
Alternatives (IIA): for any two candidate moves \(m_{1}\) and \(m_{2}\)
in a choice set \(\mathcal{N}\), the relative probability
\(P\left( m_{1} \right)/P\left( m_{2} \right)\) is the same regardless
of which other candidates are in \(\mathcal{N}\). In the reassignment
context, candidate moves are not exchangeable: a move that targets an
amino acid already serviced by many codons is plausibly more
substitutable with similar moves than the IIA structure allows.

We adopt IIA here because the goal is \emph{explanatory} rather than
\emph{predictive}: the test asks whether topology adds explanatory value
beyond physicochemical cost (LR test, \(\Delta\)AICc) within the same
candidate set, not whether the model accurately predicts which specific
reassignment will occur next. IIA remains an \textbf{untested structural
assumption} of the estimator, however, even for the explanatory reading:
every candidate contributes to the choice-set denominator, so the
composition and similarity of the non-observed alternatives affect the
fitted likelihood, coefficients, and LR statistics. The
restricted-candidate refits in \Cref{sec:s-iia-restricted} (and the
alternative-specification fits reported alongside them) are best read as
\textbf{sensitivity analyses} against candidate-set composition, not as
evidence that IIA is harmless. A mixed logit relaxing IIA would be more
appropriate for both explanatory calibration and prediction; future work
could pursue this once a larger event set permits identification of
mixed-logit covariance parameters. The parametric predictive simulation
reported in the main text (§3.5) provides a calibration check on the
model's marginal feature distribution rather than just in-sample AICc:
the event-step topology-breaking rate under \(Q_{6}\) adjacency with the
\textbf{increase-in-components} (\(\Delta\beta_{0} > 0\)) definition,
observed at 5/66 = 0.076, is reproduced by the M3 simulation at mean
\(0.077 \pm 0.033\) (\(n =\)10,000 draws; parametric predictive
\(p = 0.60\)). This is the event-step / \(Q_{6}\) /
\(\Delta\beta_{0} > 0\) rate; the lineage-collapsed / \(H(3,4)\) /
new-disconnection rate quoted for the primary topology-avoidance test in
§3.4 (6/28 = 0.21) is a different quantity and is not directly
comparable.

A second, structurally distinct concern about the candidate set is its
\textbf{composition}: whether the universe of \(\approx\) 1,280 moves
contains so many biologically implausible alternatives that the fitted
topology coefficient is partly identifying their absence rather than a
topology-specific avoidance effect. We address this concern in
\Cref{sec:s-iia-restricted} by refitting M1--M4 on candidate sets that
have been pruned to biologically plausible alternatives only.

\subsection{Candidate-set composition: restricted-candidate
sensitivity}\label{sec:s-iia-restricted}

A separate concern about the candidate set is that the universe of
\(\approx\)1,280 single-codon moves admits biologically catastrophic
alternatives (reassigning AUG-Met, simultaneous multi-codon changes
implicit in the single-step framing, reassignments to stop in essential
codons) that natural selection has already removed from the option set.
Models with strongly negative coefficients on \(\Delta_{\text{topo}}\)
and \(\Delta_{\text{phys}}\) may therefore be partly rediscovering that
natural reassignments are not biologically catastrophic, inflating ΔAICc
magnitudes beyond what the explanatory thesis (topology adds value
beyond physicochemistry) strictly requires. The qualitative claim is
unaffected by this concern, but the magnitudes need calibration.

We address this by refitting M1--M4 (and the \(H(3,4)\) verification
variants) on a \textbf{restricted candidate set}: at each event-step we
retain only candidates whose target amino acid is already serviced by a
codon at Hamming distance \(\leq d\) from the reassigned codon (i.e.,
\(\Delta_{\text{tRNA }} \leq d\)), with
\(d \in \left\{ 1,2,3 \right\}\). The observed move is always retained
regardless of its \(\Delta_{\text{tRNA}}\) so the likelihood remains
well-defined. We report all three filters as a \textbf{bracketing}
sensitivity rather than designating any single \(d\) as the calibrated
biological cut, because the relevant decoding literature does not pin a
single distance. The strictest cut, \(d = 1\), is the closest match to
canonical wobble-position mismatch tolerance \citep{crick1966} and is
the most defensible biological-plausibility floor; \(d = 3\) permits
near-cognate routes and is included as a loose upper bound. We had
previously labelled the intermediate \(d = 2\) filter as ``primary'' in
earlier drafts; in the present version we report it as the midpoint of
the bracket, since the qualitative reading of the table does not depend
on that label and the reviewers correctly noted that designating
\(d = 2\) ``primary'' without an independent biological anchor risks
reporting the most interpretable effect size rather than the most
defensible one. \Cref{tbl:s-condlogit-restricted} reports the resulting
\(\Delta\)AICc gaps under all four candidate sets.

\begin{longtable}[]{@{}
  >{\centering\arraybackslash}p{(\linewidth - 12\tabcolsep) * \real{0.1429}}
  >{\raggedleft\arraybackslash}p{(\linewidth - 12\tabcolsep) * \real{0.1429}}
  >{\raggedleft\arraybackslash}p{(\linewidth - 12\tabcolsep) * \real{0.1429}}
  >{\raggedleft\arraybackslash}p{(\linewidth - 12\tabcolsep) * \real{0.1429}}
  >{\raggedleft\arraybackslash}p{(\linewidth - 12\tabcolsep) * \real{0.1429}}
  >{\raggedleft\arraybackslash}p{(\linewidth - 12\tabcolsep) * \real{0.1429}}
  >{\raggedleft\arraybackslash}p{(\linewidth - 12\tabcolsep) * \real{0.1429}}@{}}
\caption{Restricted-candidate sensitivity for the conditional-logit
comparison. Each row refits M1, M2, M3, M4 (and the \(H(3,4)\) topology
variants) on a candidate set filtered to
\(\Delta_{\text{tRNA }} \leq d\), where \(\Delta_{\text{tRNA}}\) is the
Hamming distance from the reassigned codon to the nearest existing codon
for the target amino acid. Observed moves are always retained so
likelihoods remain comparable. The ``Unrestricted'' row reproduces the
unrestricted main-text numbers. ΔAICc gaps shrink as the candidate set
is restricted. This is expected: removing biologically implausible
candidates removes contrasts the model otherwise exploits, so the table
should be read as a \textbf{bracket}: the strictest
biological-plausibility floor at \(d = 1\) (~275 candidates) gives
ΔAICc(M1\(\rightarrow\)M3) \(\approx 14\), just above the conventional
Burnham--Anderson reference, while the looser \(d = 2\) filter (~727
candidates) gives \(\approx 60\) and the looser \(d = 3\) filter
\(\approx 95\). ΔAICc(M2\(\rightarrow\)M3) stays large at every
threshold. The qualitative claim ``topology adds explanatory value
beyond physicochemistry'' is therefore robust to candidate-set
composition; the \textbf{magnitude} of the topology--physicochemistry
separation is best characterised by the \(d = 1\) floor rather than by
the unrestricted \(\approx 110\) figure. The ΔAICc(M3\(\rightarrow\)M4)
column under the restricted filters is \textbf{not} interpretable,
because the filter is defined on \(\Delta_{\text{tRNA}}\) so the M4
tRNA-feature distribution shifts mechanically with the filter; the
unrestricted-set ΔAICc(M3\(\rightarrow\)M4) of ~2 is the calibrated
reading and shows the heuristic tRNA proxy is
uninformative.}\tabularnewline
\toprule\noalign{}
\textbf{Filter} & \textbf{Mean cand.} & \textbf{Obs. kept} &
\textbf{\vtop{\hbox{\strut \(\Delta\)AICc}\hbox{\strut M1\(\rightarrow\)M3}}}
&
\textbf{\vtop{\hbox{\strut \(\Delta\)AICc}\hbox{\strut M2\(\rightarrow\)M3}}}
&
\textbf{\vtop{\hbox{\strut \(\Delta\)AICc}\hbox{\strut M3\(\rightarrow\)M4}}}
&
\textbf{\vtop{\hbox{\strut \(\Delta\)AICc}\hbox{\strut M1\(\rightarrow\)M3\textsubscript{\(H(3,4)\)}}}} \\
\midrule\noalign{}
\endfirsthead
\toprule\noalign{}
\textbf{Filter} & \textbf{Mean cand.} & \textbf{Obs. kept} &
\textbf{\vtop{\hbox{\strut \(\Delta\)AICc}\hbox{\strut M1\(\rightarrow\)M3}}}
&
\textbf{\vtop{\hbox{\strut \(\Delta\)AICc}\hbox{\strut M2\(\rightarrow\)M3}}}
&
\textbf{\vtop{\hbox{\strut \(\Delta\)AICc}\hbox{\strut M3\(\rightarrow\)M4}}}
&
\textbf{\vtop{\hbox{\strut \(\Delta\)AICc}\hbox{\strut M1\(\rightarrow\)M3\textsubscript{\(H(3,4)\)}}}} \\
\midrule\noalign{}
\endhead
\bottomrule\noalign{}
\endlastfoot
Unrestricted & 1280 & --- & 108 & 89 & -2.1 & 91 \\
\(\Delta_{\text{tRNA }} \leq 3\) & 1088 & 66 / 66 & 95 & 85 & -1.6 &
77 \\
\(\Delta_{\text{tRNA }} \leq 2\) & 727 & 66 / 66 & 60 & 77 & 17.3 &
41 \\
\(\Delta_{\text{tRNA }} \leq 1\) & 275 & 66 / 66 & 14 & 73 & 95.3 &
15 \\
\end{longtable}

\protect\phantomsection\label{tbl:s-condlogit-restricted}{}

Two structural points are worth flagging beyond the table itself. First,
the ΔAICc(M1\(\rightarrow\)M3) gap is exactly the quantity sensitive to
candidate-set composition: physicochemistry alone (M1) is a flat
baseline that benefits most from the implausible-candidate filler being
present, so removing it shrinks the M1\(\rightarrow\)M3 distance more
than the M2\(\rightarrow\)M3 distance. Second,
ΔAICc(M2\(\rightarrow\)M3) is robust across all three thresholds
(\ensuremath{\approx}73--85), reflecting that physicochemistry contributes the same
incremental signal regardless of which candidates surround the observed
events. The overall reading is therefore conservative: the
unrestricted-set ΔAICc(M1\(\rightarrow\)M3) \ensuremath{\approx} 110 over-states the
\textbf{magnitude} of the topology--physicochemistry separation but does
not over-state its existence, and the \(d = 1\) floor (\(\approx 14\))
provides the most defensible biological-plausibility-bounded effect size
and shows the qualitative claim survives even under the strictest cut.

\section{Conditional logit: clade-exclusion
sensitivity}\label{sec:s-condlogit-clade}

The conditional-logit framework can confound a single clade's
reassignment events with a global topology-avoidance signal if that
clade is unusually extensive (or unusually idiosyncratic). To bound this
concern, we refit M1--M4 under each of seven clade-exclusion regimes
(\Cref{tbl:s-condlogit-clade}), removing the indicated NCBI translation
tables from the event-step list and re-running the full
conditional-logit pipeline end-to-end (build candidate sets
\(\rightarrow\) enumerate event orderings up to \(k!\) for \(k \leq 6\)
\(\rightarrow\) vectorised maximum-likelihood fit \(\rightarrow\) ΔAICc
against M1 and M2). The seven exclusion regimes match those applied to
the hypergeometric topology-avoidance test (\Cref{sec:s-clade}), which
themselves follow the phylogenetic-distribution analysis in
\citep{sengupta2007}.

\begin{longtable}[]{@{}
  >{\raggedright\arraybackslash}p{(\linewidth - 10\tabcolsep) * \real{0.1667}}
  >{\raggedright\arraybackslash}p{(\linewidth - 10\tabcolsep) * \real{0.1667}}
  >{\centering\arraybackslash}p{(\linewidth - 10\tabcolsep) * \real{0.1667}}
  >{\raggedleft\arraybackslash}p{(\linewidth - 10\tabcolsep) * \real{0.1667}}
  >{\raggedleft\arraybackslash}p{(\linewidth - 10\tabcolsep) * \real{0.1667}}
  >{\raggedleft\arraybackslash}p{(\linewidth - 10\tabcolsep) * \real{0.1667}}@{}}
\caption{\textbf{Per-regime conditional-logit clade-exclusion
sensitivity (7 regimes).} Each row refits M1--M4 with the indicated NCBI
translation tables removed; \(n\) = remaining event-steps. M3
(physicochemistry + \(Q_{6}\) topology) remains favored over both M1
(physicochemistry only) and M2 (topology only) in every regime. The
minimum \(\Delta\)AICc(M1\(\rightarrow\)M3) across all seven regimes is
49; the maximum is 117. The conventional Burnham--Anderson reference of
\(\Delta\)AICc\(> 10\) was calibrated on linear-regression contexts, so
we use it as a reference threshold rather than a formally calibrated
cut-off in this conditional-logit setting; the parametric predictive
simulation reported in main-text §3.5 (event-step topology-breaking rate
under \(Q_{6}\) / \(\Delta\beta_{0} > 0\): observed 5/66 = 0.076 vs
simulated 0.077, \(p = 0.60\)) is the more directly interpretable
calibration check. \(\Delta\)AICc(M3\(\rightarrow\)M4) values near 2
indicate that adding the heuristic tRNA-distance proxy provides no
incremental fit.}\tabularnewline
\toprule\noalign{}
\begin{minipage}[b]{\linewidth}\raggedright
\textbf{Excluded clade}
\end{minipage} & \begin{minipage}[b]{\linewidth}\raggedright
\textbf{Tables out}
\end{minipage} & \begin{minipage}[b]{\linewidth}\centering
\textbf{\(n\)}
\end{minipage} & \begin{minipage}[b]{\linewidth}\raggedleft
\textbf{\(\Delta\)AICc\\
M1\(\rightarrow\)M3}\strut
\end{minipage} & \begin{minipage}[b]{\linewidth}\raggedleft
\textbf{\(\Delta\)AICc\\
M2\(\rightarrow\)M3}\strut
\end{minipage} & \begin{minipage}[b]{\linewidth}\raggedleft
\textbf{\(\Delta\)AICc\\
M3\(\rightarrow\)M4}\strut
\end{minipage} \\
\midrule\noalign{}
\endfirsthead
\toprule\noalign{}
\begin{minipage}[b]{\linewidth}\raggedright
\textbf{Excluded clade}
\end{minipage} & \begin{minipage}[b]{\linewidth}\raggedright
\textbf{Tables out}
\end{minipage} & \begin{minipage}[b]{\linewidth}\centering
\textbf{\(n\)}
\end{minipage} & \begin{minipage}[b]{\linewidth}\raggedleft
\textbf{\(\Delta\)AICc\\
M1\(\rightarrow\)M3}\strut
\end{minipage} & \begin{minipage}[b]{\linewidth}\raggedleft
\textbf{\(\Delta\)AICc\\
M2\(\rightarrow\)M3}\strut
\end{minipage} & \begin{minipage}[b]{\linewidth}\raggedleft
\textbf{\(\Delta\)AICc\\
M3\(\rightarrow\)M4}\strut
\end{minipage} \\
\midrule\noalign{}
\endhead
\bottomrule\noalign{}
\endlastfoot
ciliates & 6, 10, 15, 27, 28, 29, 30 & 52 & 83.7 & 31.6 & 1.8 \\
yeast mito & 3 & 60 & 103.1 & 87.7 & 2.2 \\
CUG clade & 12, 26 & 64 & 116.7 & 94.9 & 2.2 \\
metazoan mito & 2, 5, 9, 13, 14, 21 & 40 & 48.8 & 88.6 & 0.4 \\
algal mito & 16, 22 & 63 & 115.5 & 85.7 & 2 \\
protist mito & 4, 23 & 63 & 100.7 & 92.4 & 1.9 \\
hemichordate mito & 24, 33 & 59 & 92.3 & 83.5 & 2.1 \\
\end{longtable}

\protect\phantomsection\label{tbl:s-condlogit-clade}{}

\section{Per-table optimality: standard-code-proximity
audit}\label{sec:s-proximity}

This audit grounds the disaggregation reported in main-text §3.3, where
we separated the per-table BH-FDR result into ``informative-distance''
tables (\(\geq 3\) reassignments from standard) and ``near-standard''
tables (\(\leq 2\) reassignments).

The methodological concern: a variant code that differs from the
standard by only a few reassignments has a per-table quartet-pattern
shuffle null distribution dominated by permutations very close to the
standard code, simply because few permutations of the variant's block
structure yield codes that are far from standard. In that limit, ``table
\(X\) falls in the bottom 5\% of permutations preserving \(X\)`s block
structure'' partly tests whether \(X\) is close to the standard code
(which it is, by construction), not whether \(X\) is independently
optimal.

For each NCBI translation table we computed three quantities alongside
the unconditional per-table quantile: (i) the Hamming distance \(d_{H}\)
from the standard code (number of codons with a different AA label) for
both the observed variant code and each quartet-pattern shuffle null
draw, (ii) the variant's null quantile \textbf{conditional} on null
draws within \(\pm 2\) codons of the variant's \(d_{H}\), and (iii) the
fraction of null draws whose \(d_{H}\) is at most the variant's
\(d_{H}\) (a complementary ``proximity rank'' of the variant against the
null). The motivation for the conditional quantile was to ask whether
variants with low unconditional \(p\)-values would still appear
unusually low-cost when restricted to null draws of equivalent
\(d_{H}\). As \Cref{tbl:s-proximity} shows, the conditional bucket
turns out to be empty for every variant: under block-preserving
permutation, every null draw has \(d_{H} \geq 30\) from the standard
code (the lowest observed \(d_{H}\) in the null distribution across all
27 tables is 30, while the highest variant \(d_{H}\) in the registry is
6), so no null draws fall within \(\pm 2\) codons of any variant's
\(d_{H}\). The conditional quantile is therefore \textbf{not
informative} about whether a variant's optimality is intrinsic versus
standard-code-proximity-driven.

The proximity-rank diagnostic in the rightmost column is informative:
every variant in the registry has \(d_{H}\) smaller than every null
draw, so all variants sit at the extreme low-\(d_{H}\) tail of the
quartet-pattern shuffle null distribution. This means each variant's
per-table \(p\)-value is, structurally, partly a
proximity-to-standard-code measurement, and the ``informative-distance''
vs ``near-standard'' disaggregation in main-text §3.3 is therefore based
on absolute \(d_{H}\) thresholds (3 reassignments) rather than on this
audit's conditional quantile.

\begin{longtable}[]{@{}>{\centering\arraybackslash}p{(\linewidth - 10\tabcolsep) * \real{0.08}}>{\centering\arraybackslash}p{(\linewidth - 10\tabcolsep) * \real{0.12}}>{\raggedleft\arraybackslash}p{(\linewidth - 10\tabcolsep) * \real{0.15}}>{\raggedleft\arraybackslash}p{(\linewidth - 10\tabcolsep) * \real{0.15}}>{\raggedleft\arraybackslash}p{(\linewidth - 10\tabcolsep) * \real{0.14}}>{\raggedleft\arraybackslash}p{(\linewidth - 10\tabcolsep) * \real{0.36}}@{}}
\caption{\textbf{Per-table standard-code-proximity audit.} For each NCBI
translation table other than the standard code (table 1), \(d_{H}\) is
the number of codons with a different AA label from the standard code.
\textbf{Null \(d_{H}\) range} is the (min, max) of \(d_{H}\) across the
10,000 quartet-pattern shuffle null draws for that table.
\textbf{Quantile (uncond.)} is the variant's per-table
block-preserving-null quantile from main-text §3.3 (lower = more
error-minimising). \textbf{\(n\) in \(d_{H} \pm 2\) bucket} counts null
draws with \(d_{H}\) within \(\pm 2\) of the variant's \(d_{H}\) (would
have been the conditional-quantile denominator); this is zero for every
table because every quartet-pattern shuffle null draw has
\(d_{H} \geq 30\) while every variant has \(d_{H} \leq 6\).
\textbf{Frac. null \(d_{H} \leq d_{H}^{\text{obs}}\)} is the fraction of
null draws closer to (or as close as) the variant in \(d_{H}\); this is
uniformly \(0\)\%, confirming that every variant sits at the extreme
low-\(d_{H}\) tail of its quartet-pattern shuffle null. The
conditional-quantile diagnostic is therefore not informative on this
data, and the disaggregation in main-text §3.3 rests on the absolute
\(d_{H} \geq 3\) threshold rather than on this conditional
analysis.}\tabularnewline
\toprule\noalign{}
\textbf{Table} & \textbf{\(d_{H}\) (variant)} & \textbf{Null \(d_{H}\)
range} & \textbf{Quantile (uncond.)} & \textbf{\(n\) in \(d_{H} \pm 2\)
bucket} & \textbf{Frac. null \(d_{H} \leq d_{H}^{\text{obs}}\)} \\
\midrule\noalign{}
\endfirsthead
\toprule\noalign{}
\textbf{Table} & \textbf{\(d_{H}\) (variant)} & \textbf{Null \(d_{H}\)
range} & \textbf{Quantile (uncond.)} & \textbf{\(n\) in \(d_{H} \pm 2\)
bucket} & \textbf{Frac. null \(d_{H} \leq d_{H}^{\text{obs}}\)} \\
\midrule\noalign{}
\endhead
\bottomrule\noalign{}
\endlastfoot
\multicolumn{6}{@{}l@{}}{%
\textbf{Informative-distance tables (\(d_{H} \geq 3\) codon
reassignments from standard).}} \\
3 & 6 & 37--60 & 7.4\% & 0 & 0.0\% \\
21 & 5 & 39--60 & 1.4\% & 0 & 0.0\% \\
14 & 5 & 38--61 & 2.2\% & 0 & 0.0\% \\
33 & 4 & 41--61 & 2.0\% & 0 & 0.0\% \\
13 & 4 & 32--60 & 2.4\% & 0 & 0.0\% \\
9 & 4 & 36--60 & 1.4\% & 0 & 0.0\% \\
5 & 4 & 38--60 & 2.4\% & 0 & 0.0\% \\
2 & 4 & 36--58 & 1.4\% & 0 & 0.0\% \\
31 & 3 & 44--64 & 0.3\% & 0 & 0.0\% \\
28 & 3 & 44--64 & 0.2\% & 0 & 0.0\% \\
27 & 3 & 37--64 & 0.2\% & 0 & 0.0\% \\
24 & 3 & 40--60 & 1.5\% & 0 & 0.0\% \\
\multicolumn{6}{@{}l@{}}{%
\textbf{Near-standard tables (\(d_{H} \leq 2\) codon reassignments from
standard).}} \\
30 & 2 & 38--60 & 0.0\% & 0 & 0.0\% \\
29 & 2 & 38--60 & 2.4\% & 0 & 0.0\% \\
23 & 2 & 35--57 & 0.8\% & 0 & 0.0\% \\
22 & 2 & 30--54 & 0.8\% & 0 & 0.0\% \\
6 & 2 & 40--60 & 0.3\% & 0 & 0.0\% \\
32 & 1 & 35--57 & 1.5\% & 0 & 0.0\% \\
26 & 1 & 32--56 & 1.0\% & 0 & 0.0\% \\
25 & 1 & 36--60 & 1.1\% & 0 & 0.0\% \\
16 & 1 & 33--57 & 1.9\% & 0 & 0.0\% \\
15 & 1 & 35--57 & 0.3\% & 0 & 0.0\% \\
12 & 1 & 35--56 & 0.8\% & 0 & 0.0\% \\
10 & 1 & 40--60 & 0.5\% & 0 & 0.0\% \\
4 & 1 & 38--60 & 1.0\% & 0 & 0.0\% \\
11 & 0 & 34--56 & 1.2\% & 0 & 0.0\% \\
\end{longtable}

\protect\phantomsection\label{tbl:s-proximity}{}

The full per-table audit, including \(d_{H}\) summary statistics for the
null distribution and the per-table observed score, is released in
\texttt{output/coloring\_optimality.json} under the
\texttt{per\_table\_proximity\_audit} key.

\begin{figure}
\centering
\emfig{FigB\_per\_table\_optimality}
\caption{Coloring optimality preserved across variant codes. Per-table
Grantham quartet-pattern shuffle null quantile for each of the 27 NCBI
translation tables, tested against its own quartet-pattern shuffle null
(\(n =\)10,000 per table). Bars sorted by quantile; the dashed reference
line marks the 5\% threshold. Companion to main-text §3.3 and Figure 2C;
the ``informative-distance'' vs ``near-standard'' disaggregation of §3.3
is not visible in this plot because it is a proximity-audit output
rather than a null-quantile statistic (see \Cref{tbl:s-proximity}).}
\end{figure}

\protect\phantomsection\label{fig:s-per-table}{}

\section{Topology avoidance: clade-exclusion sensitivity
(hypergeometric/permutation)}\label{sec:s-clade}

This section provides the hypergeometric/permutation counterpart to the
conditional-logit clade-exclusion analysis in
\Cref{sec:s-condlogit-clade}. The two tests probe related but distinct
quantities: \Cref{sec:s-condlogit-clade} tests whether the
\(\Delta_{\text{topo}}\) regression coefficient survives clade
exclusion; this section tests whether the global landscape-vs-observed
depletion ratio survives clade exclusion. Following the
phylogenetic-distribution analysis of mitochondrial reassignment events
by \citep{sengupta2007}, we iteratively excluded each of seven major
taxonomic groups (\Cref{tbl:s-phylo-clade}) and re-ran the \(Q_{6}\)
topology-avoidance hypergeometric on the reduced event set, holding the
1,280-move candidate landscape fixed.

\begin{longtable}[]{@{}
  >{\raggedright\arraybackslash}p{(\linewidth - 10\tabcolsep) * \real{0.1667}}
  >{\raggedright\arraybackslash}p{(\linewidth - 10\tabcolsep) * \real{0.1667}}
  >{\centering\arraybackslash}p{(\linewidth - 10\tabcolsep) * \real{0.1667}}
  >{\raggedleft\arraybackslash}p{(\linewidth - 10\tabcolsep) * \real{0.1667}}
  >{\raggedleft\arraybackslash}p{(\linewidth - 10\tabcolsep) * \real{0.1667}}
  >{\raggedleft\arraybackslash}p{(\linewidth - 10\tabcolsep) * \real{0.1667}}@{}}
\caption{\textbf{Hypergeometric clade-exclusion sensitivity for the
\(Q_{6}\) topology-avoidance test.} Each row removes the indicated NCBI
translation tables (matching the clade definitions of
\citep{sengupta2007}) from the de-duplicated event list and re-runs the
hypergeometric depletion test against the same 1,280-move candidate
landscape. \textbf{Breakers} = topology-breaking events (new
disconnection, \(Q_{6}\) adjacency); \textbf{Rate} = observed breaker
rate (vs candidate-landscape rate of 72.7\% under the default encoding).
All seven exclusions yield \(p < 10^{- 5}\). Excluding yeast
mitochondrial (NCBI translation table 3) \textbf{strengthens} the
depletion to \(p \approx 3.6 \times 10^{- 11}\) because that table
contributes 4 of the 6 lineage-collapsed topology-breakers; this is the
denominator-effect noted in main-text §3.4.}\tabularnewline
\toprule\noalign{}
\begin{minipage}[b]{\linewidth}\raggedright
\textbf{Excluded clade}
\end{minipage} & \begin{minipage}[b]{\linewidth}\raggedright
\textbf{Tables out}
\end{minipage} & \begin{minipage}[b]{\linewidth}\centering
\textbf{\(n\)}
\end{minipage} & \begin{minipage}[b]{\linewidth}\raggedleft
\textbf{Breakers}
\end{minipage} & \begin{minipage}[b]{\linewidth}\raggedleft
\textbf{Rate (\%)}
\end{minipage} & \begin{minipage}[b]{\linewidth}\raggedleft
\textbf{Hyper. \(p\)}
\end{minipage} \\
\midrule\noalign{}
\endfirsthead
\toprule\noalign{}
\begin{minipage}[b]{\linewidth}\raggedright
\textbf{Excluded clade}
\end{minipage} & \begin{minipage}[b]{\linewidth}\raggedright
\textbf{Tables out}
\end{minipage} & \begin{minipage}[b]{\linewidth}\centering
\textbf{\(n\)}
\end{minipage} & \begin{minipage}[b]{\linewidth}\raggedleft
\textbf{Breakers}
\end{minipage} & \begin{minipage}[b]{\linewidth}\raggedleft
\textbf{Rate (\%)}
\end{minipage} & \begin{minipage}[b]{\linewidth}\raggedleft
\textbf{Hyper. \(p\)}
\end{minipage} \\
\midrule\noalign{}
\endhead
\bottomrule\noalign{}
\endlastfoot
ciliates & 6, 10, 15, 27, 28, 29, 30 & 24 & 6 & 25 &
\(1.22 \times 10^{- 6}\) \\
yeast mito & 3 & 24 & 2 & 8.3 & \(3.64 \times 10^{- 11}\) \\
CUG clade & 12, 26 & 26 & 5 & 19.2 & \(1.41 \times 10^{- 8}\) \\
metazoan mito & 2, 5, 9, 13, 14, 21 & 22 & 6 & 27.3 &
\(9.78 \times 10^{- 6}\) \\
algal mito & 16, 22 & 26 & 5 & 19.2 & \(1.41 \times 10^{- 8}\) \\
protist mito & 4, 23 & 27 & 6 & 22.2 & \(4.78 \times 10^{- 8}\) \\
hemichordate mito & 24, 33 & 27 & 6 & 22.2 & \(4.78 \times 10^{- 8}\) \\
\end{longtable}

\protect\phantomsection\label{tbl:s-phylo-clade}{}

In every exclusion, the \(Q_{6}\) depletion remains highly significant
(\(p < 10^{- 5}\)), confirming that the \(\approx 3.4\)-fold depletion
of topology-breaking changes is a pan-taxonomic pattern, not an artifact
of any single lineage.

The primary-cell \(H(3,4)\) / \(\Delta\beta_{0} > 0\) analogue of the
same seven-row exclusion is \Cref{tbl:s-phylo-clade-k43}. It uses the
encoding-independent nucleotide-adjacency graph and the
increase-in-components definition matched to the conditional-logit
feature (main-text §3.5), and confirms that the \(\approx 3.1\)-fold
\(H(3,4)\) depletion is also robust to clade exclusion.

\begin{longtable}[]{@{}
  >{\raggedright\arraybackslash}p{(\linewidth - 10\tabcolsep) * \real{0.1667}}
  >{\raggedright\arraybackslash}p{(\linewidth - 10\tabcolsep) * \real{0.1667}}
  >{\centering\arraybackslash}p{(\linewidth - 10\tabcolsep) * \real{0.1667}}
  >{\raggedleft\arraybackslash}p{(\linewidth - 10\tabcolsep) * \real{0.1667}}
  >{\raggedleft\arraybackslash}p{(\linewidth - 10\tabcolsep) * \real{0.1667}}
  >{\raggedleft\arraybackslash}p{(\linewidth - 10\tabcolsep) * \real{0.1667}}@{}}
\caption{\textbf{Hypergeometric clade-exclusion sensitivity for the
primary \(H(3,4)\) topology-avoidance test.} Each row removes the
indicated NCBI translation tables (matching the clade definitions of
\citep{sengupta2007}) from the de-duplicated event list and re-runs the
hypergeometric depletion test under the encoding-independent
\(H(3,4) = K_{4}^{3}\) adjacency with the increase-in-components
(\(\Delta\beta_{0} > 0\)) topology-breaking definition (candidate
landscape: 846 breakers of 1,280 moves = 66.1\%). \textbf{Breakers} =
topology-breaking events by \(\Delta\beta_{0} > 0\); \textbf{Rate} =
observed breaker rate. All seven exclusions yield \(p < 10^{- 3}\).
Excluding yeast mitochondrial (NCBI translation table 3) strengthens the
depletion to \(p \approx 4.2 \times 10^{- 9}\) because that table
contributes 4 of the 6 lineage-collapsed \(H(3,4)\)-breakers --- the
direct calculation quoted in main-text §3.4 for the yeast-mito-excluded
regime.}\tabularnewline
\toprule\noalign{}
\begin{minipage}[b]{\linewidth}\raggedright
\textbf{Excluded clade}
\end{minipage} & \begin{minipage}[b]{\linewidth}\raggedright
\textbf{Tables out}
\end{minipage} & \begin{minipage}[b]{\linewidth}\centering
\textbf{\(n\)}
\end{minipage} & \begin{minipage}[b]{\linewidth}\raggedleft
\textbf{Breakers}
\end{minipage} & \begin{minipage}[b]{\linewidth}\raggedleft
\textbf{Rate (\%)}
\end{minipage} & \begin{minipage}[b]{\linewidth}\raggedleft
\textbf{Hyper. \(p\)}
\end{minipage} \\
\midrule\noalign{}
\endfirsthead
\toprule\noalign{}
\begin{minipage}[b]{\linewidth}\raggedright
\textbf{Excluded clade}
\end{minipage} & \begin{minipage}[b]{\linewidth}\raggedright
\textbf{Tables out}
\end{minipage} & \begin{minipage}[b]{\linewidth}\centering
\textbf{\(n\)}
\end{minipage} & \begin{minipage}[b]{\linewidth}\raggedleft
\textbf{Breakers}
\end{minipage} & \begin{minipage}[b]{\linewidth}\raggedleft
\textbf{Rate (\%)}
\end{minipage} & \begin{minipage}[b]{\linewidth}\raggedleft
\textbf{Hyper. \(p\)}
\end{minipage} \\
\midrule\noalign{}
\endhead
\bottomrule\noalign{}
\endlastfoot
ciliates & 6, 10, 15, 27, 28, 29, 30 & 24 & 6 & 25 &
\(3.96 \times 10^{- 5}\) \\
yeast mito & 3 & 24 & 2 & 8.3 & \(4.21 \times 10^{- 9}\) \\
CUG clade & 12, 26 & 26 & 4 & 15.4 & \(1.08 \times 10^{- 7}\) \\
metazoan mito & 2, 5, 9, 13, 14, 21 & 22 & 6 & 27.3 &
\(2.03 \times 10^{- 4}\) \\
algal mito & 16, 22 & 26 & 6 & 23.1 & \(7.29 \times 10^{- 6}\) \\
protist mito & 4, 23 & 27 & 6 & 22.2 & \(3.07 \times 10^{- 6}\) \\
hemichordate mito & 24, 33 & 27 & 6 & 22.2 & \(3.07 \times 10^{- 6}\) \\
\end{longtable}

\protect\phantomsection\label{tbl:s-phylo-clade-k43}{}

The clade-exclusion robustness is a denominator effect (removing clades
that contributed zero or one topology-breaking event each leaves the
breakage rate essentially unchanged) rather than a numerator effect; the
\textbf{avoidance} of topology-breaking moves is what propagates across
many lineages, while the \textbf{attempts} to break topology are
concentrated in yeast mitochondrial. Detailed per-clade counts are
released as \texttt{output/phylogenetic\_sensitivity.json} (the
\(Q_{6}\) / new-disconnection cell) and
\texttt{output/phylogenetic\_sensitivity\_k43.json} (the \(H(3,4)\) /
\(\Delta\beta_{0} > 0\) cell).

\section{Complete tRNA gene count data}\label{sec:s-trna}

This section documents the empirical foundation for the tRNA enrichment
analysis in main-text §3.6. We ran tRNAscan-SE 2.0.12 on 18 NCBI genome
assemblies; 15 of these 18 organisms enter the 24-pairing
Fisher--Stouffer enrichment analysis as either variant-code or
standard-code-control repertoires (\Cref{tbl:s-trna-provenance} below).
The remaining 3 (\emph{Blastocrithidia nonstop} P57;
\emph{Mycoplasmoides genitalium} G37; \emph{Mycoplasmoides pneumoniae}
M129) are used only as \textbf{mechanistic boundary cases} in main-text
§4.3: their reassignment routes (anticodon stem shortening and anticodon
modification respectively) act on a single tRNA gene without changing
the gene-count distribution, so a Fisher-count enrichment test would
report ``no enrichment'' and misinterpret the mechanism as an
absence-of-response. Adding literature/database counts for a small
number of legacy repertoires (see \Cref{tbl:s-trna-provenance}
``source-class'' column: annotation, GtRNAdb, literature) brings the
total pairing set to 24 pairings across 24 organisms with mixed
provenance.

The Fisher-exact denominator convention used throughout is the
\textbf{by-amino-acid sum}: for a focal amino acid \(a\) in variant
genome \(V\) the \(2 \times 2\) table is
\(\left\lbrack \left\lbrack a_{V},\left( \sum_{a' \neq a}a'_{V} \right) \right\rbrack,\left\lbrack a_{C},\left( \sum_{a' \neq a}a'_{C} \right) \right\rbrack \right\rbrack\),
where the by-amino-acid sum is over the \textbf{Std20} column of
\Cref{tbl:s-trna}. This excludes SeC, undetermined-isotype, suppressor,
and pseudogene rows so that the ratio \(\frac{a}{\sum_{a'}}a'\) has a
consistent interpretation across genomes (fractional tRNA-repertoire
share for the focal AA). The 24 exact \(2 \times 2\) tables and
per-pairing Fisher \(p\)-values are printed in
\Cref{tbl:s-trna-provenance} and the machine-readable CSV
\texttt{output/tables/T-trna-24row-provenance.csv}. The first-pass
Isotype/Anticodon totals reported parenthetically in the Reassigned-AA
column of \Cref{tbl:s-trna} below are provided so the reader can see
the per-anticodon breakdown; they are \textbf{not} the Fisher
denominators.

\subsection{\texorpdfstring{24-pairing input table with provenance and
Fisher \(2 \times 2\)
counts}{24-pairing input table with provenance and Fisher 2 \textbackslash times 2 counts}}\label{sec:s-trna-provenance}

\begin{longtable}[]{@{}llclccrrrr@{}}
\caption{\textbf{Complete 24-pairing tRNA-enrichment input table.} For
each pairing: variant and control organism (short key; full organism
name in \Cref{tbl:s-trna}), reassigned amino acid, brief reassignment
description, whether the reassignment creates a new \(Q_{6}\) AA-family
disconnection at \(\varepsilon = 1\) (Y/n), and the \(H(3,4)\) impact
(creates / extends / no-effect / n/a for pairings without a \(Q_{6}\)
disconnection). \textbf{Var \(\frac{a}{N}\)}, \textbf{Ctl
\(\frac{a}{N}\)}: focal AA count over the by-amino-acid-sum denominator.
\textbf{OR}: Fisher-exact odds ratio (variant vs control).
\textbf{Fisher \(p\)}: one-sided upper-tail \(p\)-value. Source-class
breakdown across the 48 organism-slots (24 variant + 24 control):
tRNAscan-SE 2.0.12 = 38, literature = 4, annotation = 4, GtRNAdb = 2.
Unique organisms: 24; unique organisms with tRNAscan-SE-verified counts:
15. Machine-readable CSV:
\texttt{output/tables/T-trna-24row-provenance.csv}.}\tabularnewline
\toprule\noalign{}
\textbf{Variant} & \textbf{Control} & \textbf{AA} &
\textbf{Reassignment} & \textbf{\(Q_{6}\) break?} & \textbf{H(3,4)} &
\textbf{Var \(\frac{a}{N}\)} & \textbf{Ctl \(\frac{a}{N}\)} &
\textbf{OR} & \textbf{Fisher \(p\)} \\
\midrule\noalign{}
\endfirsthead
\toprule\noalign{}
\textbf{Variant} & \textbf{Control} & \textbf{AA} &
\textbf{Reassignment} & \textbf{\(Q_{6}\) break?} & \textbf{H(3,4)} &
\textbf{Var \(\frac{a}{N}\)} & \textbf{Ctl \(\frac{a}{N}\)} &
\textbf{OR} & \textbf{Fisher \(p\)} \\
\midrule\noalign{}
\endhead
\bottomrule\noalign{}
\endlastfoot
scerevisiae & ylipolytica & Thr & CUA: Leu\ensuremath{\to}{}Thr; CUC: Leu\ensuremath{\to}{}Thr; C\ldots{}
& Y & creates & 2/24 & 1/27 & 2.36 & 0.455 \\
sobliquus & creinhardtii & Leu & UAG: Stop\ensuremath{\to}{}Leu & Y & no-effect & 2/22 &
1/16 & 1.5 & 0.621 \\
ptannophilus & lthermotolerans & Ala & CUG: Leu\ensuremath{\to}{}Ala & Y & creates &
14/207 & 11/206 & 1.29 & 0.345 \\
calbicans & lthermotolerans & Ser & CUG: Leu\ensuremath{\to}{}Ser & Y & extends & 16/207
& 13/206 & 1.24 & 0.355 \\
tthermophila & imultifiliis & Gln & UAA: Stop\ensuremath{\to}{}Gln; UAG: Stop\ensuremath{\to}{}Gln & n &
no-effect & 54/712 & 3/141 & 3.78 & 0.008 \\
ptetraurelia & imultifiliis & Gln & UAA: Stop\ensuremath{\to}{}Gln; UAG: Stop\ensuremath{\to}{}Gln & n &
no-effect & 18/215 & 3/141 & 4.2 & 0.01 \\
tthermophila & scoeruleus & Gln & UAA: Stop\ensuremath{\to}{}Gln; UAG: Stop\ensuremath{\to}{}Gln & n &
no-effect & 54/712 & 11/267 & 1.91 & 0.032 \\
ptetraurelia & scoeruleus & Gln & UAA: Stop\ensuremath{\to}{}Gln; UAG: Stop\ensuremath{\to}{}Gln & n &
no-effect & 18/215 & 11/267 & 2.13 & 0.04 \\
oxytricha & scoeruleus & Gln & UAA: Stop\ensuremath{\to}{}Gln; UAG: Stop\ensuremath{\to}{}Gln & n &
no-effect & 8/90 & 11/267 & 2.27 & 0.075 \\
eoctocarinatus & scoeruleus & Cys & UGA: Stop\ensuremath{\to}{}Cys & n & no-effect &
5/114 & 7/267 & 1.7 & 0.272 \\
bjaponicum & scoeruleus & Trp & (no direct reassignment for t\ldots{} &
n & no-effect & 5/115 & 6/267 & 1.98 & 0.21 \\
eaediculatus & scoeruleus & Cys & UGA: Stop\ensuremath{\to}{}Cys & n & no-effect & 4/77 &
7/267 & 2.04 & 0.215 \\
eaediculatus & fsalina & Cys & UGA: Stop\ensuremath{\to}{}Cys & n & no-effect & 4/77 &
1/86 & 4.66 & 0.151 \\
eamieti & scoeruleus & Cys & UGA: Stop\ensuremath{\to}{}Cys & n & no-effect & 8/116 &
7/267 & 2.75 & 0.049 \\
efocardii & fsalina & Cys & UGA: Stop\ensuremath{\to}{}Cys & n & no-effect & 3/59 & 1/86
& 4.55 & 0.184 \\
eparawoodruffi & scoeruleus & Cys & UGA: Stop\ensuremath{\to}{}Cys & n & no-effect &
9/145 & 7/267 & 2.46 & 0.065 \\
eweissei & fsalina & Cys & UGA: Stop\ensuremath{\to}{}Cys & n & no-effect & 20/463 & 1/86
& 3.84 & 0.132 \\
ewoodruffi & scoeruleus & Cys & UGA: Stop\ensuremath{\to}{}Cys & n & no-effect & 4/79 &
7/267 & 1.98 & 0.226 \\
ppersalinus & imultifiliis & Gln & UAA: Stop\ensuremath{\to}{}Gln; UAG: Stop\ensuremath{\to}{}Gln & n &
no-effect & 20/260 & 3/141 & 3.83 & 0.015 \\
ppersalinus & fsalina & Gln & UAA: Stop\ensuremath{\to}{}Gln; UAG: Stop\ensuremath{\to}{}Gln & n &
no-effect & 20/260 & 3/86 & 2.31 & 0.132 \\
hgrandinella & scoeruleus & Gln & UAA: Stop\ensuremath{\to}{}Gln; UAG: Stop\ensuremath{\to}{}Gln & n &
no-effect & 9/129 & 11/267 & 1.75 & 0.165 \\
hgrandinella & fsalina & Gln & UAA: Stop\ensuremath{\to}{}Gln; UAG: Stop\ensuremath{\to}{}Gln & n &
no-effect & 9/129 & 3/86 & 2.08 & 0.218 \\
bstoltei & scoeruleus & Trp & (no direct reassignment for t\ldots{} & n
& no-effect & 6/166 & 6/267 & 1.63 & 0.289 \\
bstoltei & fsalina & Trp & (no direct reassignment for t\ldots{} & n &
no-effect & 6/166 & 2/86 & 1.57 & 0.447 \\
\end{longtable}

\protect\phantomsection\label{tbl:s-trna-provenance}{}

\subsection{tRNAscan-SE verified organisms}

All tRNA gene counts were obtained by running tRNAscan-SE 2.0.12
\citep{chan2019} with Infernal 1.1.4 on NCBI genome assemblies.
Eukaryotic organisms were scanned in \texttt{-E} mode;
\emph{Mycoplasmoides} species in \texttt{-B} (bacterial) mode. Total
counts are Infernal-confirmed (the more conservative count after the
second-pass filter) and match the Total column of \Cref{tbl:s-trna}
below; the summary that underlies the main-text Fisher--Stouffer
combination is reported as main-text Table 7. The five-column breakdown
in \Cref{tbl:s-trna} sums exactly to Total: Std20 (decoding the
standard 20 amino acids) + SeC (selenocysteine, anticodon UCA, scanned
via the dedicated \texttt{TRNAinf-euk-SeC.cm} model) + Supp (possible
suppressor tRNAs with CTA/TTA/UCA anticodons) + Undet (predicted tRNAs
whose isotype could not be determined) + Pseudo (predicted pseudogenes
filtered by the Infernal second-pass isotype validation).

\begin{longtable}[]{@{}
  >{\raggedright\arraybackslash}p{(\linewidth - 18\tabcolsep) * \real{0.1200}}
  >{\centering\arraybackslash}p{(\linewidth - 18\tabcolsep) * \real{0.0400}}
  >{\raggedright\arraybackslash}p{(\linewidth - 18\tabcolsep) * \real{0.2000}}
  >{\raggedleft\arraybackslash}p{(\linewidth - 18\tabcolsep) * \real{0.0600}}
  >{\raggedleft\arraybackslash}p{(\linewidth - 18\tabcolsep) * \real{0.0600}}
  >{\raggedleft\arraybackslash}p{(\linewidth - 18\tabcolsep) * \real{0.0500}}
  >{\raggedleft\arraybackslash}p{(\linewidth - 18\tabcolsep) * \real{0.0600}}
  >{\raggedleft\arraybackslash}p{(\linewidth - 18\tabcolsep) * \real{0.0600}}
  >{\raggedleft\arraybackslash}p{(\linewidth - 18\tabcolsep) * \real{0.0700}}
  >{\raggedright\arraybackslash}p{(\linewidth - 18\tabcolsep) * \real{0.2800}}@{}}
\caption{Complete tRNAscan-SE 2.0.12 results. The \textbf{Supp} column
and the \textbf{Std20/SeC/Undet/Pseudo} columns are Infernal-confirmed
second-pass counts and sum exactly to \textbf{Total}: \textbf{Std20} =
standard 20-AA tRNAs; \textbf{SeC} = selenocysteine (anticodon UCA);
\textbf{Supp} = possible suppressor tRNAs (CTA/TTA/UCA anticodons that
could read stop codons) \textbf{after} Infernal's isotype-validation
filter; \textbf{Undet} = predicted tRNAs whose isotype could not be
determined; \textbf{Pseudo} = predicted pseudogenes filtered by the
second-pass validation. The Reassigned-AA breakdown, in contrast, uses
tRNAscan-SE's \textbf{first-pass} Isotype/Anticodon counts, which are
the per-anticodon totals reported directly by the covariance-model scan
before the Infernal second-pass reclassification. Consequently the
first-pass suppressor total in the Reassigned-AA parenthetical can
exceed the second-pass \textbf{Supp} column when Infernal reclassifies
borderline hits into \textbf{Std20} or \textbf{Pseudo}; the differences
here are 0--2 tRNAs per organism (e.g. \emph{T. thermophila}: 37
second-pass suppressors versus 39 first-pass = 9 CTA + 30 TTA; \emph{E.
parawoodruffi}: 3 versus 5 = 0 CTA + 1 TTA + 4 UCA). All Fisher-exact
tests in Table 7 and §S10 use the first-pass per-amino-acid totals shown
in the Reassigned-AA column, which are the biologically interpretable
per-anticodon counts. \textsuperscript{§}\emph{Blepharisma stoltei} is
indexed as NCBI translation table 15 (Blepharisma Nuclear Code, defined
as UAG\(\rightarrow\)Gln) by virtue of its genus assignment, but the
strain-specific MAC genome reads UGA\(\rightarrow\)Trp via a dedicated
suppressor tRNA-Trp(UCA) \citep{singh2023}. The analysis here tests
Trp-tRNA enrichment on that empirical reading rather than the legacy
table-15 nominal code. \textsuperscript{*}Anticodon stem shortening
(4-bp stem rather than canonical 5-bp).
\textsuperscript{†}Post-transcriptional modification (a single
tRNA-Trp(CCA) reads both UGG and UGA after base
modification).}\tabularnewline
\toprule\noalign{}
\begin{minipage}[b]{\linewidth}\raggedright
\textbf{Organism}
\end{minipage} & \begin{minipage}[b]{\linewidth}\centering
\textbf{Tbl}
\end{minipage} & \begin{minipage}[b]{\linewidth}\raggedright
\textbf{Assembly}
\end{minipage} & \begin{minipage}[b]{\linewidth}\raggedleft
\textbf{Total}
\end{minipage} & \begin{minipage}[b]{\linewidth}\raggedleft
\textbf{Std20}
\end{minipage} & \begin{minipage}[b]{\linewidth}\raggedleft
\textbf{SeC}
\end{minipage} & \begin{minipage}[b]{\linewidth}\raggedleft
\textbf{Supp}
\end{minipage} & \begin{minipage}[b]{\linewidth}\raggedleft
\textbf{Undet}
\end{minipage} & \begin{minipage}[b]{\linewidth}\raggedleft
\textbf{Pseudo}
\end{minipage} & \begin{minipage}[b]{\linewidth}\raggedright
\textbf{Reassigned AA}
\end{minipage} \\
\midrule\noalign{}
\endfirsthead
\toprule\noalign{}
\begin{minipage}[b]{\linewidth}\raggedright
\textbf{Organism}
\end{minipage} & \begin{minipage}[b]{\linewidth}\centering
\textbf{Tbl}
\end{minipage} & \begin{minipage}[b]{\linewidth}\raggedright
\textbf{Assembly}
\end{minipage} & \begin{minipage}[b]{\linewidth}\raggedleft
\textbf{Total}
\end{minipage} & \begin{minipage}[b]{\linewidth}\raggedleft
\textbf{Std20}
\end{minipage} & \begin{minipage}[b]{\linewidth}\raggedleft
\textbf{SeC}
\end{minipage} & \begin{minipage}[b]{\linewidth}\raggedleft
\textbf{Supp}
\end{minipage} & \begin{minipage}[b]{\linewidth}\raggedleft
\textbf{Undet}
\end{minipage} & \begin{minipage}[b]{\linewidth}\raggedleft
\textbf{Pseudo}
\end{minipage} & \begin{minipage}[b]{\linewidth}\raggedright
\textbf{Reassigned AA}
\end{minipage} \\
\midrule\noalign{}
\endhead
\bottomrule\noalign{}
\endlastfoot
\emph{T. thermophila} & 6 & GCF\_000189635.1 & 718 & 672 & 1 & 37 & 4 &
4 & Gln: 54 (15+39) \\
\emph{P. tetraurelia} & 6 & GCF\_000165425.1 & 216 & 202 & 1 & 11 & 0 &
2 & Gln: 18 (7+11) \\
\emph{O. trifallax} & 6 & GCA\_000295675.1 & 94 & 83 & 2 & 6 & 2 & 1 &
Gln: 8 (2+6) \\
\emph{P. persalinus} & 6 & GCA\_001447515.1 & 262 & 228 & 1 & 15 & 0 &
18 & Gln: 20 (5+15) \\
\emph{H. grandinella} & 6 & GCA\_006369765.1 & 130 & 121 & 1 & 3 & 0 & 5
& Gln: 9 (6+3) \\
\emph{E. aediculatus} & 10 & GCA\_030463445.1 & 80 & 76 & 1 & 2 & 1 & 0
& Cys: 4 (3+1 UCA) \\
\emph{E. amieti} & 10 & GCA\_048569255.1 & 120 & 103 & 1 & 6 & 1 & 9 &
Cys: 8 (4+4 UCA) \\
\emph{E. focardii} & 10 & GCA\_001880345.2 & 62 & 56 & 1 & 3 & 0 & 2 &
Cys: 3 (1+2 UCA) \\
\emph{E. parawoodruffi} & 10 & GCA\_021440025.1 & 149 & 128 & 0 & 3 & 3
& 15 & Cys: 9 (5+4 UCA) \\
\emph{E. weissei} & 10 & GCA\_021440005.1 & 495 & 390 & 0 & 12 & 14 & 79
& Cys: 20 (17+3 UCA) \\
\emph{E. woodruffi} & 10 & GCA\_027382605.1 & 83 & 74 & 1 & 3 & 1 & 4 &
Cys: 4 (2+2 UCA) \\
\emph{B. stoltei}\textsuperscript{§} & 15 & GCA\_965603825.1 & 169 & 165
& 1 & 0 & 1 & 2 & Trp: 6 \\
\emph{B. nonstop} P57 & 31 & GCA\_028554745.1 & 68 & 65 & 1 & 2 & 0 & 0
& Trp: 2\textsuperscript{*} \\
\emph{M. genitalium} & 4 & GCA\_000027325.1 & 36 & 35 & 0 & 1 & 0 & 0 &
Trp: 1\textsuperscript{†} \\
\emph{M. pneumoniae} & 4 & GCF\_910574535.1 & 37 & 36 & 0 & 1 & 0 & 0 &
Trp: 1\textsuperscript{†} \\
\emph{S. coeruleus} & 1 & GCA\_001970955.1 & 272 & 265 & 1 & 1 & 3 & 2 &
--- \\
\emph{I. multifiliis} & 1 & GCF\_000220395.1 & 150 & 141 & 1 & 8 & 0 & 0
& --- \\
\emph{F. salina} & 1 & GCA\_022984795.1 & 89 & 85 & 1 & 1 & 1 & 1 &
--- \\
\end{longtable}

\protect\phantomsection\label{tbl:s-trna}{}

\subsection{MIS (maximal independent set) analysis}

To address non-independence from shared control organisms, we
constructed a conflict graph in which edges connect pairings sharing an
organism. Enumerating all maximal independent sets (MIS) is equivalent
to enumerating all maximal cliques of the complement graph, which we
perform via the Bron--Kerbosch algorithm applied to the complement (with
pivoting). Each MIS represents a set of pairings where no two share an
organism and no additional pairing can be added without creating a
conflict.

The 24-pairing conflict graph admits \textbf{332} MIS, all of size 6.
Across these 332 sets the Stouffer combined \(p\)-value distribution is
as follows: median \(p = 0.037\), best \(p = 0.012\), worst
\(p = 0.104\); 264 of 332 (79.5\%) fall below the 0.05 threshold and 0
of 332 below 0.01. The extreme cases:

\begin{itemize}
\item
  \textbf{Best-case MIS} (\(p = 0.012\)): \emph{S. cerevisiae} mito/Thr,
  \emph{S. obliquus} mito/Leu, \emph{P. tannophilus}/Ala, \emph{P.
  tetraurelia}/Gln, \emph{T. thermophila}/Gln, \emph{P. persalinus}/Gln.
\item
  \textbf{Worst-case MIS} (\(p = 0.104\)): \emph{S. cerevisiae}
  mito/Thr, \emph{S. obliquus} mito/Leu, \emph{C. albicans}/Ser,
  \emph{E. octocarinatus}/Cys, \emph{P. persalinus}/Gln, \emph{B.
  stoltei}/Trp.
\end{itemize}

Because the median just clears the 0.05 threshold while the worst case
does not, the tRNA enrichment result is treated as \textbf{exploratory}
in the claim hierarchy (\Cref{tbl:s-claims}): the signal is present in
the median-independent-subset sense but is not robust to worst-case
subset choice.

\textbf{Independence-graph caveat.} The conflict graph encodes shared
\textbf{organisms}, not phylogenetic distance. Two ciliates from
different orders are treated as independent in this construction even
though they may share more evolutionary signal than two arbitrary
eukaryotes; at the same time, every MIS retains the \emph{S. cerevisiae}
mito/Thr pairing whose tRNA counts come from GtRNAdb rather than from a
tRNAscan-SE 2.0.12 run on the assembly (see §S10.3). A
phylogenetic-distance--based independence criterion would yield a
smaller effective \(n\) (the same 4--6 phylogenetically independent
origins discussed in main-text Limitations); the MIS analysis should be
read as bounding subset variability under the
organism-shared-by-pairings criterion, not as a phylogeny-aware test.

\textbf{Topology-breaking subset.} As a pre-specified mechanistic
complement to the all-pairings analysis, we restrict the
Fisher--Stouffer combination to the four pairings whose underlying
reassignment creates a new amino-acid disconnection under \(Q_{6}\):
\emph{S. cerevisiae} mito/Thr (table 3), \emph{S. obliquus} mito/Leu
(table 22, equivalent to chlorophycean mito table 16), \emph{P.
tannophilus}/Ala (table 26), and \emph{C. albicans}/Ser (table 12). This
subset combination yields Stouffer \(p = 0.387\), i.e. null. The
all-pairings signal is therefore driven largely by topology-preserving
stop-to-sense reassignment systems (UAR\(\rightarrow\)Gln in ciliates,
UGA\(\rightarrow\)Cys in \emph{Euplotes}, UGA\(\rightarrow\)Trp in
\emph{Blepharisma}), and the tRNA-duplication mechanism is not evidenced
by the topology-breaking cases in isolation. This is why the claim
hierarchy treats the tRNA result as exploratory-decoding-accommodation
rather than as evidence of compensation-for-disconnection.

\subsection{\texorpdfstring{\emph{Saccharomyces cerevisiae}
literature-derived
control}{Saccharomyces cerevisiae literature-derived control}}

The \emph{Saccharomyces cerevisiae} tRNA gene counts used in the
yeast-mito Thr disconnection pairing come from GtRNAdb \citep{chan2016}
rather than tRNAscan-SE 2.0.12 in this work. We retain them to keep the
well-characterized yeast-mito Thr case in the pairing set and flag the
difference in source explicitly.

\begin{figure}
\centering
\emfig{FigE\_trna\_aa\_rank}
\caption{Rank of the reassigned amino acid among all tRNA gene counts,
per variant-code \(\times\) control pairing. Rank 1 = most enriched
amino acid in the variant-code-vs-control comparison. (Most of the
pairings here are topology-preserving stop-to-sense reassignments ---
UAR\(\rightarrow\)Gln, UGA\(\rightarrow\)Cys, UGA\(\rightarrow\)Trp ---
rather than topology-breaking; the ``disconnection'' framing that
appears in earlier drafts is inaccurate for the majority of pairings.)
The distribution is left-heavy toward rank 1--3 across pairings,
consistent with heterogeneous accommodation-of-decoding across
variant-code lineages. The associated statistical summary (main-text
§3.6, Table 7): 332 MIS of size 6, median Stouffer \(p = 0.037\), worst
\(p = 0.104\); topology-breaking-only subset (\(n = 4\)) is null
(Stouffer \(p = 0.387\)).}
\end{figure}

\protect\phantomsection\label{fig:s-trna-aa-rank}{}

\section{Complete reassignment database}\label{sec:s-reassignment}

The reassignment database underpins both the topology-avoidance test
(main-text §3.4) and the conditional-logit fits (main-text §3.5). It
comprises every codon reassignment event across the 27 NCBI translation
tables, relative to the standard code (table 1). Each event records the
codon, source amino acid, target amino acid, and Hamming distance to the
nearest codon already encoding the target amino acid in the standard
code. De-duplication to unique (codon, target amino acid) pairs yields
the event set used in the topology-avoidance hypergeometric and
table-preserving permutation tests. The conditional-logit model operates
on the full per-table event-step list, which preserves recurrent
reassignments across independent lineages. Exact counts depend on the
most recent NCBI gc.prt revision (currently v4.6 retrieved 2026-04-25);
a current pipeline run yields 66 raw events de-duplicating to 28 unique
(codon, target-AA) pairs.

\textbf{Why two counts (\(n =\)28 hypergeometric vs \(n =\)66
conditional-logit).} The hypergeometric topology-avoidance test
(main-text §3.4) uses the de-duplicated \(n =\)28 unique (codon,
target-AA) pairs, since the test is a sample-from-a-finite-landscape
contrast that should not double-count the same move. The
conditional-logit comparison (main-text §3.5) uses the full per-table
\(n =\)66 event-step list because each event-step is a \textbf{separate
choice}: a recurrent reassignment such as UGA\(\rightarrow\)Trp seen
across many independent mitochondrial lineages contributes one
independent choice per lineage to the conditional likelihood, since each
lineage's reassignment is a distinct evolutionary event conditional on
that lineage's candidate set. The conditional-logit clade-exclusion
sensitivity (\Cref{sec:s-condlogit-clade}) inherits this convention:
removing a clade removes all of its event-steps from the per-table list,
including recurrent UGA\(\rightarrow\)Trp instances within that clade,
so the ΔAICc magnitudes report the \textbf{clade-removed-as-event-steps}
sensitivity. We did not additionally re-weight by recurrence (e.g., one
event per (codon, target) pair rather than per lineage); that
alternative weighting would correspond more closely to the de-duplicated
hypergeometric framing and would shrink ΔAICc magnitudes toward the
restricted-candidate \(d = 1\) floor reported in
\Cref{sec:s-iia-restricted}, but does not change the qualitative
ranking M3 \(>\) M1, M2.

\begin{longtable}[]{@{}
  >{\centering\arraybackslash}p{(\linewidth - 4\tabcolsep) * \real{0.3333}}
  >{\raggedright\arraybackslash}p{(\linewidth - 4\tabcolsep) * \real{0.3333}}
  >{\centering\arraybackslash}p{(\linewidth - 4\tabcolsep) * \real{0.3333}}@{}}
\caption{\textbf{Per-table reassignment-event counts (raw, before
de-duplication).} The 25 NCBI translation tables that differ from the
standard code account for 66 raw codon-reassignment events. Tables 1
(Standard) and 11 (Bacterial / Archaeal / Plant Plastid) share identical
sense-codon mappings and contribute zero reassignment events. Tables 27
(Karyorelict Nuclear) and 28 (Condylostoma Nuclear) likewise share
identical sense-codon mappings and are listed separately to match NCBI
numbering. The 25-table counts sum to the all-events total, which
de-duplicates to 28 unique (codon, target-AA) pairs across the registry.
The complete event-level database (with codon, source-AA, target-AA,
Hamming-to-nearest-target columns) is released as
\texttt{output/tables/T10\_reassignment\_db.csv} in the
\texttt{codontopo} repository.}\tabularnewline
\toprule\noalign{}
\begin{minipage}[b]{\linewidth}\centering
\textbf{Table}
\end{minipage} & \begin{minipage}[b]{\linewidth}\raggedright
\textbf{NCBI translation table}
\end{minipage} & \begin{minipage}[b]{\linewidth}\centering
\textbf{Events}
\end{minipage} \\
\midrule\noalign{}
\endfirsthead
\toprule\noalign{}
\begin{minipage}[b]{\linewidth}\centering
\textbf{Table}
\end{minipage} & \begin{minipage}[b]{\linewidth}\raggedright
\textbf{NCBI translation table}
\end{minipage} & \begin{minipage}[b]{\linewidth}\centering
\textbf{Events}
\end{minipage} \\
\midrule\noalign{}
\endhead
\bottomrule\noalign{}
\endlastfoot
2 & Vertebrate Mitochondrial & 4 \\
3 & Yeast Mitochondrial & 6 \\
4 & Mold/Protozoan/Coelenterate Mito & 1 \\
5 & Invertebrate Mitochondrial & 4 \\
6 & Ciliate Nuclear & 2 \\
9 & Echinoderm/Flatworm Mito & 4 \\
10 & Euplotid Nuclear & 1 \\
12 & Alternative Yeast Nuclear & 1 \\
13 & Ascidian Mitochondrial & 4 \\
14 & Alternative Flatworm Mito & 5 \\
15 & Blepharisma Macronuclear & 1 \\
16 & Chlorophycean Mito & 1 \\
21 & Trematode Mitochondrial & 5 \\
22 & Scenedesmus obliquus Mito & 2 \\
23 & Thraustochytrium Mito & 2 \\
24 & Rhabdopleuridae Mito & 3 \\
25 & Candidate Division SR1 / Gracilibacteria & 1 \\
26 & Pachysolen tannophilus Nuclear & 1 \\
27 & Karyorelict Nuclear & 3 \\
28 & Condylostoma Nuclear & 3 \\
29 & Mesodinium Nuclear & 2 \\
30 & Peritrich Nuclear & 2 \\
31 & Blastocrithidia Nuclear & 3 \\
32 & Balanophoraceae Plastid & 1 \\
33 & Cephalodiscidae Mito & 4 \\
\end{longtable}

\protect\phantomsection\label{tbl:s-reassignment}{}

\section{Structural-preservation index
(visualization-only)}\label{sec:s-feasibility}

For Figure 5A of the main text we use a heuristic
\textbf{structural-preservation index} \(S(m) \in \lbrack 0,1\rbrack\)
for each candidate single-codon reassignment \(m\) from the standard
code. The index is \textbf{not} used in any inferential test in the
manuscript; it is a visualization aid for delineating high- versus
low-structure-preserving regions of the 1,280-move candidate landscape.
We report the exact implementation here for completeness.

The index is a weighted sum of four discrete indicators. Given the
reassignment \(m\), let
\(I_{\text{tw}(m)},I_{\text{ff}(m)} \in \left\{ 0,1 \right\}\) record
whether the post-reassignment code preserves the two-fold bit-5
filtration and the four-fold prefix filtration described in main-text
§2.2; let \(I_{\text{ser}(m)} \in \left\{ 0,1 \right\}\) record whether
Serine remains disconnected at \(\varepsilon = 1\); and let
\(n_{\text{disc}(m)}\) denote the number of amino acids disconnected at
\(\varepsilon = 1\) under \(m\). Then

\[S(m) = 0.25 \cdot I_{\text{tw}(m)} + 0.25 \cdot I_{\text{ff}(m)} + 0.30 \cdot I_{\text{ser}(m)} + w\left( n_{\text{disc}(m)} \right),\]

where \(w(1) = 0.20\) (a single disconnected amino acid, as in the
standard code), \(w(0) = 0.10\) (partial credit for a fully connected
variant), and \(w(n) = 0.00\) for \(n \geq 2\). Because every component
is discrete, \(S(m)\) can take only a small number of values across the
1,280-move candidate universe; the four bars visible in Figure 5A
(\(S = 0.55,0.75,0.80,1.00\); 193/27/738/322 events) enumerate the
values actually realised. The weights are a hand-tuned prioritisation of
the four structural features and are not derived from any calibration;
the index is descriptive, not inferential. Exact implementation:
\texttt{src/codon\_topo/analysis/synbio\_feasibility.py} in the
\texttt{codontopo} repository.

\section{KRAS--Fano clinical prediction: detailed
results}\label{sec:s-kras}

The KRAS--Fano clinical prediction is the most concrete biomedical
extrapolation considered for the \({\text{GF}(2)}^{6}\) framework, and
we tested it directly to delimit scope. The conjecture: XOR (``Fano'')
relationships in \({\text{GF}(2)}^{6}\) predict enrichment of specific
amino acids at KRAS G12 co-mutation sites. For each of the six common
somatic KRAS G12 variants, we computed the predicted Fano partner codon
as \(GGU \oplus m_{i} = \text{ partner}_{i}\) where \(m_{i}\) is the
mutant codon for that variant and \(\oplus\) is bitwise XOR in
\({\text{GF}(2)}^{6}\) (\Cref{tbl:s-kras}). We then tested for
co-mutation enrichment of the predicted partner amino acid against 1,670
KRAS-mutant samples from the MSK-IMPACT pan-cancer dataset
\citep{zehir2017} using Fisher's exact test with Bonferroni correction
across the six variants.

\begin{longtable}[]{@{}>{\centering\arraybackslash}p{(\linewidth - 10\tabcolsep) * \real{0.10}}>{\centering\arraybackslash}p{(\linewidth - 10\tabcolsep) * \real{0.10}}>{\centering\arraybackslash}p{(\linewidth - 10\tabcolsep) * \real{0.14}}>{\centering\arraybackslash}p{(\linewidth - 10\tabcolsep) * \real{0.22}}>{\centering\arraybackslash}p{(\linewidth - 10\tabcolsep) * \real{0.22}}>{\raggedleft\arraybackslash}p{(\linewidth - 10\tabcolsep) * \real{0.22}}@{}}
\caption{\textbf{KRAS--Fano clinical prediction tested against
MSK-IMPACT.} For each KRAS G12 variant, the Fano partner codon is
computed as the bitwise XOR (in \({\text{GF}(2)}^{6}\)) of the wild-type
and mutant codons against a fixed reference; the predicted partner amino
acid is the standard-code translation of that codon. Fisher's exact test
for co-mutation enrichment of the predicted partner amino acid in
MSK-IMPACT samples bearing each KRAS G12 variant (\(n =\)1,670
KRAS-mutant tumours; \citep{zehir2017}) yields \(p = 1.0\) for all six
variants after Bonferroni correction across the six tests. Odds ratios
are near 1.0 throughout. The conjecture is cleanly
falsified.}\tabularnewline
\toprule\noalign{}
\textbf{Variant} & \textbf{WT codon} & \textbf{Mutant codon} &
\textbf{Fano partner codon} & \textbf{Predicted partner AA} &
\textbf{Fisher \(p\) (Bonf.)} \\
\midrule\noalign{}
\endfirsthead
\toprule\noalign{}
\textbf{Variant} & \textbf{WT codon} & \textbf{Mutant codon} &
\textbf{Fano partner codon} & \textbf{Predicted partner AA} &
\textbf{Fisher \(p\) (Bonf.)} \\
\midrule\noalign{}
\endhead
\bottomrule\noalign{}
\endlastfoot
G12V & GGU & GUU & CAC & His (H) & 1.0 \\
G12D & GGU & GAU & CUC & Leu (L) & 1.0 \\
G12A & GGU & GCU & CGC & Arg (R) & 1.0 \\
G12R & GGU & CGU & GCC & Ala (A) & 1.0 \\
G12C & GGU & UGU & ACC & Thr (T) & 1.0 \\
G12S & GGU & AGU & UCC & Ser (S) & 1.0 \\
\end{longtable}

\protect\phantomsection\label{tbl:s-kras}{}

This result cleanly separates code-level error-minimization, which
operates on the amino acid assignment structure, from mutation-level
algebraic predictions, which would require DNA polymerase errors to
respect the binary encoding, which is a biologically implausible
mechanism. The Fano-line construction was a meaningful test of whether
the binary representation has predictive power at the \textbf{mutation}
level (in addition to the \textbf{assignment} level where it does carry
signal); the answer is no, and we record this falsification explicitly
in \Cref{tbl:s-claims}.

\section{Serine minimum inter-family distance across
encodings}\label{sec:s-serine-dist}

The claim that Serine's minimum inter-family Hamming distance (UCN--AGY)
equals 4 under all 24 base-to-bit encodings is false. Of the 24
encodings, 16 yield minimum distance 2 and only 8 yield distance 4. The
distance-4 result obtains only when both nucleotide pairs distinguishing
UCN from AGY (\(U \sim A\) and \(C \sim G\) in the first two positions)
are encoded at maximal Hamming distance. The correct encoding-invariant
statement is: Serine is disconnected at \(\varepsilon = 1\) under every
encoding, and its inter-family distance (\(\geq 2\)) is the largest
among the three 6-codon amino acids (Leucine and Arginine both have
inter-family distance 1). The rejection of the distance-4 invariant is
the reason the manuscript now presents \(H(3,4)\) rather than \(Q_{6}\)
as the primary adjacency for encoding-independent claims.

\section{PSL(2,7) symmetry: pre-rejection by irrep
dimensions}\label{sec:s-psl}

The claim that PSL(2,7) is the fundamental symmetry group of the genetic
code was pre-rejected by \citep{antoneli2011}, who showed that PSL(2,7)
has no 64-dimensional irreducible representation (its irreducible
representations have dimensions 1, 3, 6, 7, and 8). A group acting on
the 64-codon set must therefore decompose that action into
representations that sum to 64. Any PSL(2,7)-based decomposition of the
codon space is a composite of the available irreps rather than a
canonical single-irrep symmetry, so PSL(2,7) does not act on the codon
set as a single symmetry group in the sense of the original claim. We
record this pre-rejection to explicitly distinguish the surviving
graph-theoretic and hypercube-coloring analyses from the falsified
algebraic-symmetry claim.

\section{Holomorphic embedding: character-identity
failure}\label{sec:s-holo}

The claim that the coordinate-wise map
\({\text{GF}(2)}^{6} \rightarrow {\mathbb{C}}^{3}\) sending base pairs
to fourth roots of unity is a holomorphic embedding extending a
character of \({\text{GF}(8)}^{\ast}\) is incorrect on two counts.
First, the domain is a finite discrete set (64 points), not a complex
manifold, so ``holomorphic'' is not defined in the standard analytic
sense. Second, treating the map \(\chi\) as a group character requires
\(\chi(x + x) = {\chi(x)}^{2}\); the map fails this since
\(i^{2} = - 1 \neq 1 = \chi(0)\). The map is a bijective coordinate
labelling but not a character or holomorphic embedding, and cannot
support the analytic properties the original claim required.

\section{Source-neighborhood burden: null result}\label{sec:s-infoneg}

The source-neighborhood burden test asks whether reassigned codons sit
in worse Hamming neighborhoods with higher Grantham distance to their
neighbors. Mann--Whitney comparison of Grantham-neighborhood sums for
reassigned versus non-reassigned codons yields \(U = 301\),
\(p = 0.70\). Reassignment is therefore not driven by local escape from
costly source neighborhoods. This null does not conflict with the
conditional-logit finding that natural events favor candidate moves with
lower \(\Delta_{\text{local}}\) (\(\beta_{\text{phys }} = - 0.004\) in
the main-text M3 fit): the two tests probe different quantities.
\(\Delta_{\text{local}}\) is the change in mismatch cost induced by a
candidate move (a destination-quality measure); the source-neighborhood
burden test asks about the absolute pre-move source-neighborhood cost.
Variant codes need not originate from unusually costly source
neighborhoods (no absolute source burden), yet still favor candidate
moves that lower local mismatch when a reassignment is triggered
(opportunistic destination selection). This indicates that the
topology-avoidance constraint operates at the global graph-connectivity
level, not at the local per-codon level.

\section{Cross-family multiple-comparison
correction}\label{sec:s-multicomp}

Correction was applied within analysis families (metric, \(\rho\)-sweep,
per-table, topology-avoidance, synthetic-recoding) rather than across
all descriptive, exploratory, and confirmatory quantities. The family
structure was the natural organisation for these analyses; we did not
file an external pre-registration. The eight test families address
conceptually distinct questions (cross-metric robustness;
\(\rho\)-interpolation between \(Q_{6}\) and \(H(3,4)\); per-NCBI-table
preservation; the \(2 \times 2\) topology-avoidance audit;
clade-exclusion sensitivity following \citep{sengupta2007} the M1--M4
discrete-choice comparison including \(H(3,4)\) variants; the 24-pairing
tRNA-enrichment Stouffer combination with MIS robustness; and the
cross-study recoding reanalysis); their test statistics are non-nested.
For transparency: under cross-family Bonferroni at
\(\alpha = \frac{0.05}{8} = 6.25 \times 10^{- 3}\), the headline
\(H(3,4)\) topology depletion (\(p = 1.3 \times 10^{- 6}\)), the
cross-metric coloring optimality (\(p \leq 0.006\) across four metrics,
already at threshold), the \(\rho\)-sweep (\(p \leq 0.006\) across five
\(\rho\)), the conditional-logit ΔAICc gap of 108 (M1\(\rightarrow\)M3,
well above any threshold), and the per-table BH--FDR result all survive.
The MIS median tRNA enrichment (\(p = 0.046\)) does \textbf{not} survive
the cross-family Bonferroni threshold (\(\alpha = 0.00625\);
\(p = 0.046\) exceeds it by more than a factor of 7), and neither does
the worst-case MIS (\(p = 0.123\)). Consistent with these outcomes, the
tRNA enrichment is classified as \textbf{exploratory} rather than
confirmatory. Note that the cross-family Bonferroni threshold applies to
a coherent set of family-level \(p\)-values; we compare the raw MIS
\(p\)-values to it as a common reference, not because the raw \(p\)s,
the within-family adjusted BH results, and the AICc gaps in this list
are exchangeable family-level statistics. A stricter hierarchical scheme
(e.g. Bonferroni across per-family summary \(p\)s, with BH within
family) would yield the same qualitative conclusions. Within-family
corrections also hold: four-metric family \(\alpha = 0.0125\),
\(\rho\)-sweep family \(\alpha = 0.01\), topology-avoidance
\(2 \times 2\) family \(\alpha = 0.0125\), Ostrov segment tests
\(\alpha = 0.0167\); the per-table tests use BH--FDR (26 of 27
significant). Exploratory analyses are labeled as such throughout the
manuscript.

\section{Event-level conditional logit model: full
details}\label{sec:s-condlogit}

This section gives the implementation-level detail for the
discrete-choice analysis reported in main-text §3.5: the candidate
universe, the three feature definitions, the optimisation routine, the
order-averaging procedure for unknown event sequences, the fitted raw
and normalised coefficients, the likelihood-ratio tests, the
predictor-confounding diagnostic, and the per-event observed-move ranks.
The complete fitted JSON output is also released as
\texttt{output/evolutionary\_simulation.json} in the public repository.

\subsection{Model specification and candidate universe}

The conditional logit model treats each observed reassignment as a
discrete choice from the set of all single-codon reassignments available
at the current code state. For a code with 64 codons and 21 possible
amino acid/stop labels, the candidate set \(\mathcal{N}(C)\) contains
exactly 1,280 moves (universe U1 in §S5: 64 codons \(\times\) 20
alternative labels, excluding identity assignments). All 1,280
candidates are evaluated at each step regardless of biological
plausibility; the model's purpose is to test whether the observed moves
are statistically distinguishable from uniform sampling given the three
feature classes. As a conditional logit, the model implies an
independence of irrelevant alternatives (IIA) structure within each
choice set; the explanatory--rather--than--predictive framing (§S6)
makes IIA tolerable for our purposes.

\subsection{Feature definitions}

\textbf{Local physicochemical mismatch change
(\(\Delta_{\text{phys}}\))}: For a reassignment of codon \(c\) from
amino acid \(a\) to \(a'\), this is the change in the sum of
\citep{grantham1974} distances across all Hamming-1 edges incident to
\(c\):

\[\Delta_{\text{phys }} = \sum_{\left\{ c,c' \right\}:d(c,c') = 1}\left\lbrack \Delta(a',\text{ code}(c')) - \Delta(a,\text{ code}(c')) \right\rbrack\]

This captures the local impact on error-minimization at the reassigned
position. When a neighboring codon is assigned Stop, we set
\(\Delta(a,\text{ Stop})\) equal to the maximum Grantham distance (215),
consistent with the stop-penalty convention used in the coloring
objective (Methods §2.2 of the main manuscript).

\textbf{Topology disruption (\(\Delta_{topo,Q\_ 6}\))}: The total
increase in connected components (at \(\varepsilon = 1\)) summed across
all amino acid codon graphs under \(Q_{6}\) adjacency. Stop codons are
excluded from the topology sum; \(\Delta_{\text{topo}}\) counts
connected-component changes only for amino acid codon families. A move
that splits one amino acid's codon family into two components
contributes \(+ 1\); a move that fragments two families contributes
\(+ 2\); topology-preserving moves contribute \(0\). The default
\(\Delta_{topo,Q\_ 6}\) uses \(Q_{6}\) adjacency (Hamming-1 in the
default \({\text{GF}(2)}^{6}\) encoding \(C \rightarrow 00\),
\(U \rightarrow 01\), \(A \rightarrow 10\), \(G \rightarrow 11\)).
Because \(Q_{6}\) adjacency is encoding-dependent (8 of 24 base-to-bit
bijections give no \(Q_{6}\) topology depletion at the
candidate-landscape level; \Cref{sec:s-encoding-sweep}), we also
compute \(\Delta_{topo,H(3,4)}\) using the encoding-independent
\(H(3,4)\) adjacency (two codons are neighbors iff they differ at
exactly one nucleotide position) and refit two model variants:
M2\textsubscript{\(H(3,4)\)} (topology-only with
\(\Delta_{topo,H(3,4)}\)) and M3\textsubscript{\(H(3,4)\)}
(physicochemistry + \(\Delta_{topo,H(3,4)}\)). The encoding-robustness
comparison \(\Delta\)AICc(M1 \(\rightarrow\)
M3\textsubscript{\(H(3,4)\)}) vs \(\Delta\)AICc(M1 \(\rightarrow\) M3)
is reported in the next subsection.

\textbf{tRNA complexity proxy (\(\Delta_{\text{tRNA}}\))}: The Hamming
distance from the reassigned codon to the nearest codon already encoding
the target amino acid in the current code. This serves as a heuristic
for the tRNA repertoire change required to service the reassigned codon,
with larger distances implying more novel tRNA machinery needed.

\subsection{Fitted coefficients}

\Cref{tbl:s-condlogit-coefs} reports the maximum-likelihood coefficient
estimates for all six model variants (M1, M2, M3, M4 under \(Q_{6}\)
topology, plus M2\textsubscript{\(H(3,4)\)} and
M3\textsubscript{\(H(3,4)\)} under encoding-independent \(H(3,4)\)
topology).

\begin{longtable}[]{@{}
  >{\raggedright\arraybackslash}p{(\linewidth - 6\tabcolsep) * \real{0.2500}}
  >{\raggedright\arraybackslash}p{(\linewidth - 6\tabcolsep) * \real{0.2500}}
  >{\centering\arraybackslash}p{(\linewidth - 6\tabcolsep) * \real{0.2500}}
  >{\centering\arraybackslash}p{(\linewidth - 6\tabcolsep) * \real{0.2500}}@{}}
\caption{Conditional logit coefficient estimates from the 27-table
re-fit. Raw coefficients are on the original feature scale; normalized
coefficients are on \(z\)-scored features (multiplying raw
\(\hat{\beta}\) by the global feature standard deviation). All
\(\hat{\beta}\) values are negative, indicating that observed
reassignment histories preferentially populate moves that reduce
physicochemical mismatch, avoid topology disruption, and (weakly) prefer
target amino acids already serviced by nearby codons. The tRNA proxy
coefficient is small and non-significant (LR \(= 0.12\), \(p = 0.73\)).
The \(H(3,4)\) verification variants (M2\textsubscript{H(3,4)},
M3\textsubscript{H(3,4)}) replace the encoding-dependent
\(\Delta_{topo,Q\_ 6}\) feature with the encoding-independent
\(\Delta_{topo,H(3,4)}\). Their fitted ΔAICc values match the main-text
encoding-robustness numbers.}\tabularnewline
\toprule\noalign{}
\begin{minipage}[b]{\linewidth}\raggedright
\textbf{Model}
\end{minipage} & \begin{minipage}[b]{\linewidth}\raggedright
\textbf{Feature}
\end{minipage} & \begin{minipage}[b]{\linewidth}\centering
\textbf{\(\hat{\beta}\) (raw)}
\end{minipage} & \begin{minipage}[b]{\linewidth}\centering
\textbf{\(\hat{\beta}\) (normalized)}
\end{minipage} \\
\midrule\noalign{}
\endfirsthead
\toprule\noalign{}
\begin{minipage}[b]{\linewidth}\raggedright
\textbf{Model}
\end{minipage} & \begin{minipage}[b]{\linewidth}\raggedright
\textbf{Feature}
\end{minipage} & \begin{minipage}[b]{\linewidth}\centering
\textbf{\(\hat{\beta}\) (raw)}
\end{minipage} & \begin{minipage}[b]{\linewidth}\centering
\textbf{\(\hat{\beta}\) (normalized)}
\end{minipage} \\
\midrule\noalign{}
\endhead
\bottomrule\noalign{}
\endlastfoot
M1 & \(\Delta_{\text{phys}}\) & -4.9e-3 & -1.33 \\
M2 (\(Q_{6}\)) & \(\Delta_{topo,Q\_ 6}\) & -3.5838 & -1.67 \\
M3 (\(Q_{6}\)) & \(\Delta_{\text{phys}}\) & -4.3e-3 & -1.17 \\
M3 (\(Q_{6}\)) & \(\Delta_{topo,Q\_ 6}\) & -3.3158 & -1.54 \\
M4 & \(\Delta_{\text{phys}}\) & -4.3e-3 & -1.17 \\
M4 & \(\Delta_{topo,Q\_ 6}\) & -2.9855 & -1.39 \\
M4 & \(\Delta_{\text{tRNA}}\) & -0.2033 & -0.21 \\
M2 (\(H(3,4)\)) & \(\Delta_{topo,H(3,4)}\) & -3.6957 & -1.79 \\
M3 (\(H(3,4)\)) & \(\Delta_{\text{phys}}\) & -4.6e-3 & -1.24 \\
M3 (\(H(3,4)\)) & \(\Delta_{topo,H(3,4)}\) & -3.452 & -1.67 \\
\end{longtable}

\protect\phantomsection\label{tbl:s-condlogit-coefs}{}

\subsection{Likelihood ratio tests}

\Cref{tbl:s-condlogit-lrt} reports nested-model likelihood-ratio
statistics under both \(Q_{6}\) and \(H(3,4)\) topology variants.

\begin{longtable}[]{@{}
  >{\raggedright\arraybackslash}p{(\linewidth - 8\tabcolsep) * \real{0.2000}}
  >{\raggedright\arraybackslash}p{(\linewidth - 8\tabcolsep) * \real{0.2000}}
  >{\centering\arraybackslash}p{(\linewidth - 8\tabcolsep) * \real{0.2000}}
  >{\centering\arraybackslash}p{(\linewidth - 8\tabcolsep) * \real{0.2000}}
  >{\centering\arraybackslash}p{(\linewidth - 8\tabcolsep) * \real{0.2000}}@{}}
\caption{Likelihood-ratio tests for nested conditional logit models,
including both \(Q_{6}\) topology variants (legacy primary) and
\(H(3,4)\) topology verification variants. Both topology features (Q\_6
added to M1, H(3,4) added to M1) and physicochemistry (added to M2 /
M2\textsubscript{H(3,4)}) provide highly significant improvements. The
tRNA-complexity proxy does not improve on the phys+topo
model.}\tabularnewline
\toprule\noalign{}
\begin{minipage}[b]{\linewidth}\raggedright
\textbf{Restricted}
\end{minipage} & \begin{minipage}[b]{\linewidth}\raggedright
\textbf{Full}
\end{minipage} & \begin{minipage}[b]{\linewidth}\centering
\textbf{LR}
\end{minipage} & \begin{minipage}[b]{\linewidth}\centering
\textbf{df}
\end{minipage} & \begin{minipage}[b]{\linewidth}\centering
\textbf{\(p\)}
\end{minipage} \\
\midrule\noalign{}
\endfirsthead
\toprule\noalign{}
\begin{minipage}[b]{\linewidth}\raggedright
\textbf{Restricted}
\end{minipage} & \begin{minipage}[b]{\linewidth}\raggedright
\textbf{Full}
\end{minipage} & \begin{minipage}[b]{\linewidth}\centering
\textbf{LR}
\end{minipage} & \begin{minipage}[b]{\linewidth}\centering
\textbf{df}
\end{minipage} & \begin{minipage}[b]{\linewidth}\centering
\textbf{\(p\)}
\end{minipage} \\
\midrule\noalign{}
\endhead
\bottomrule\noalign{}
\endlastfoot
M1 (phys) & M3 (phys+topo, \(Q_{6}\)) & 110.4 & 1 & \(\ll 10^{- 10}\) \\
M2 (topo, \(Q_{6}\)) & M3 (phys+topo, \(Q_{6}\)) & 91.2 & 1 &
\(\ll 10^{- 10}\) \\
M3 (phys+topo, \(Q_{6}\)) & M4 (full) & 0.1 & 1 & \(0.73\) \\
M1 (phys) & M3\textsubscript{\(H(3,4)\)} (phys+topo, \(H(3,4)\)) & 93.4
& 1 & \(\ll 10^{- 10}\) \\
M2\textsubscript{\(H(3,4)\)} (topo, \(H(3,4)\)) &
M3\textsubscript{\(H(3,4)\)} (phys+topo, \(H(3,4)\)) & 97.3 & 1 &
\(\ll 10^{- 10}\) \\
\end{longtable}

\protect\phantomsection\label{tbl:s-condlogit-lrt}{}

\subsection{Confounding diagnostic}

Across \(\approx\)84,000 candidate moves pooled from all choice sets,
the Spearman correlation between \(\Delta_{\text{phys}}\) and
\(\Delta_{topo,Q\_ 6}\) is \(\rho = 0.15\)
(\Cref{fig:s-condlogit-diagnostics}, panel A), indicating that the two
predictors carry largely independent information. Moves that reduce
physicochemical mismatch are only slightly more likely to also preserve
topology, and the conditional logit framework accounts for any residual
collinearity through simultaneous estimation. Coefficient estimates with
95\% CIs (panel B) and parametric predictive replication of the observed
topology-breaking rate (panel C) complete the diagnostic set.

\begin{figure}
\centering
\emfig{FigG\_condlogit\_diagnostics}
\caption{Conditional-logit model diagnostics. Fit quality and
confounding diagnostics for the six candidate models M1--M4 +
\(H(3,4)\)-verification variants (likelihood-ratio tests were applied
only to nested pairs): joint distribution of \(\Delta_{\text{phys}}\) vs
\(\Delta_{\text{topo}}\) across the candidate landscape (weakly
correlated, \(\rho = 0.15\)), coefficient estimates with 95\% confidence
intervals, and parametric predictive replication of the observed
topology-breaking rate under M3. Companion to main-text Figure 4 (which
reports the AICc comparison and observed-move percentile-rank
distribution).}
\end{figure}

\protect\phantomsection\label{fig:s-condlogit-diagnostics}{}

\subsection{Observed move percentile ranks}

Under the best model (M3), observed natural reassignments rank at the
89.5th percentile on average (mean rank 134 out of 1,280 candidates).
Notable individual events: UGA\(\rightarrow\)Trp ranks consistently
above the 98th percentile across all tables where it occurs, indicating
that this reassignment is the highest-ranked under the combined
phys+topo score. The one clear outlier is the yeast mitochondrial
CUU\(\rightarrow\)Thr reassignment (30th percentile), consistent with
NCBI translation table 3's status as the sole marginal exception in the
per-table optimality analysis.

\subsection{Order-averaging implementation}

For tables with \(k > 1\) reassignment events, the temporal ordering is
unknown. We marginalized the conditional logit likelihood over all
\(k!\) orderings, computing
\(L_{\text{table }} = (1/k!)\sum_{\sigma}\prod_{s}P\left( m_{\sigma(s)}^{\ast}~|~\mathcal{N}(C_{\sigma,s}) \right)\)
using log-sum-exp for numerical stability. For all tables with
\(k \leq 6\) events (which covers every table in the dataset, including
the largest, yeast mitochondrial, \(k = 6\), \(6! = 720\) orderings), we
enumerate all orderings exactly. The implementation precomputes feature
matrices once per (table, ordering, step) into stacked numpy arrays,
after which each Nelder--Mead / L-BFGS-B function evaluation reduces to
a small batch of matrix multiplies and \texttt{scipy.special.logsumexp}
calls; the four nested models are fit concurrently on a single shared
bundle via \texttt{joblib} threads to avoid memory duplication.

\section{ProtSub matrix and metric correlation
analysis}\label{sec:s-protsub}

This section evaluates whether the standard-code optimality result is
confined to a single amino-acid distance scale. We quantify metric
overlap across amino-acid pairs and null-code scores, and report a
structure-aware ProtSub sensitivity analysis with the caveat that
ProtSub is alignment-derived and therefore code-dependent in the sense
of the \citep{digiulio2001} critique. The analyses motivate the
treatment in Methods §2.2 and Results §3.1 of the main text.

\subsection{ProtSub as a code-dependent robustness check}

ProtSub is derived from co-evolving residue pairs in 2,320 Pfam
multiple-sequence alignments of proteins encoded under the standard
genetic code; it is therefore a code-dependent (MSA-derived)
substitution matrix in the sense of \citep{digiulio2001}. We converted
the EMBOSS log-odds form to a positive distance via the standard
diagonal-anchored relation
\(d(i,j) = \left( s(i,i) + s(j,j) \right)/2 - s(i,j)\), which guarantees
\(d(i,i) = 0\) and yields strictly positive off-diagonal entries (range
1.0 to 15.0). Sanity values: \(d(I,L) = 2.0\), \(d(K,R) = 4.0\),
\(d(F,Y) = 4.0\), \(d(D,E) = 3.5\), \(d(C,W) = 14.5\).

Under the same quartet-pattern shuffle null used for the four primary
metrics (\(n =\)10,000, seed 135325), the standard code achieves
\(F = 1089.5\) against null mean \(1180.2 \pm 29.7\) (\(z = 3.05\),
quantile \(0.04\%\), \(p =\)4 times 10\^{}(-4)\(\)). This is the most
extreme percentile of any metric in the panel. \citep{buschmann2026}
report that BLOSUM62, an AlphaFold-derived substitution matrix (AFSM),
and 16 other matrix families perform similarly across multiple
sequence-alignment tasks, with the interpretation that substitution
matrices implicitly encode physicochemical reality; their result bounds
the \citep{digiulio2001} tautology rather than eliminating it. We
therefore treat ProtSub as a robustness check rather than as a primary
test.

\subsection{Metric correlation matrices}

For each pair of metrics, we computed Spearman correlations at two
levels: (i) across the 190 unordered amino-acid pairs of the 20 standard
amino acids; and (ii) across 2,000 random codes drawn from the same
quartet-pattern shuffle null used for the optimality tests. The
AA-pair-level correlations bound how similar two metrics''
\textbf{distance scales} are; the code-level correlations bound how
similarly they \textbf{rank} codes.

\begin{longtable}[]{@{}lrrrrr@{}}
\caption{Pairwise Spearman correlations of amino-acid distance values
across the 190 unordered AA pairs (20 standard amino acids). PR = Woese
polar requirement. KD = Kyte--Doolittle hydropathy. Median off-diagonal
\(\rho = 0.49\); minimum 0.23 (KD vs ProtSub).}\tabularnewline
\toprule\noalign{}
& \textbf{Grantham} & \textbf{Miyata} & \textbf{PR} & \textbf{KD} &
\textbf{ProtSub} \\
\midrule\noalign{}
\endfirsthead
\toprule\noalign{}
& \textbf{Grantham} & \textbf{Miyata} & \textbf{PR} & \textbf{KD} &
\textbf{ProtSub} \\
\midrule\noalign{}
\endhead
\bottomrule\noalign{}
\endlastfoot
Grantham & 1.00 & 0.71 & 0.47 & 0.36 & 0.63 \\
Miyata & 0.71 & 1.00 & 0.73 & 0.54 & 0.49 \\
PR & 0.47 & 0.73 & 1.00 & 0.44 & 0.26 \\
KD & 0.36 & 0.54 & 0.44 & 1.00 & 0.23 \\
ProtSub & 0.63 & 0.49 & 0.26 & 0.23 & 1.00 \\
\end{longtable}

\protect\phantomsection\label{tbl:s-metric-corr-pairs}{}

\begin{longtable}[]{@{}lrrrrr@{}}
\caption{Pairwise Spearman correlations of \(F\)-scores across 2,000
random codes drawn from the quartet-pattern shuffle null. Higher than
the AA-pair-level correlations (\Cref{tbl:s-metric-corr-pairs}), as
expected when distances are summed over the 192 \(Q_{6}\) edges. Median
off-diagonal \(\rho = 0.82\).}\tabularnewline
\toprule\noalign{}
& \textbf{Grantham} & \textbf{Miyata} & \textbf{PR} & \textbf{KD} &
\textbf{ProtSub} \\
\midrule\noalign{}
\endfirsthead
\toprule\noalign{}
& \textbf{Grantham} & \textbf{Miyata} & \textbf{PR} & \textbf{KD} &
\textbf{ProtSub} \\
\midrule\noalign{}
\endhead
\bottomrule\noalign{}
\endlastfoot
Grantham & 1.00 & 0.85 & 0.75 & 0.71 & 0.83 \\
Miyata & 0.85 & 1.00 & 0.87 & 0.89 & 0.82 \\
PR & 0.75 & 0.87 & 1.00 & 0.78 & 0.74 \\
KD & 0.71 & 0.89 & 0.78 & 1.00 & 0.79 \\
ProtSub & 0.83 & 0.82 & 0.74 & 0.79 & 1.00 \\
\end{longtable}

\protect\phantomsection\label{tbl:s-metric-corr-codes}{}

\begin{longtable}[]{@{}lr@{}}
\caption{Partial Spearman correlations of code-level \(F\)-scores (2,000
null codes), with each pair's relationship controlled for all three
other metrics simultaneously. Five of ten pairs become approximately
uncorrelated after adjustment
(\(\left| \rho_{\text{partial}} \right| < 0.1\); no formal confidence
intervals or \(p\)-values are reported here, and ``approximately
uncorrelated'' is a descriptive statement about the point estimate
rather than a test of statistical independence), consistent with each
metric contributing unique variance not fully explained by the other
four. The high pairwise correlations in \Cref{tbl:s-metric-corr-codes}
reflect a shared physicochemical signal; the partial correlations
isolate the residual unique contribution of each metric.}\tabularnewline
\toprule\noalign{}
\textbf{Metric pair} & \textbf{Partial Spearman \(\rho\)} \\
\midrule\noalign{}
\endfirsthead
\toprule\noalign{}
\textbf{Metric pair} & \textbf{Partial Spearman \(\rho\)} \\
\midrule\noalign{}
\endhead
\bottomrule\noalign{}
\endlastfoot
Grantham \(\sim\) Miyata & +0.48 \\
Grantham \(\sim\) polar requirement & +0.03 \\
Grantham \(\sim\) Kyte--Doolittle & \ensuremath{-}0.31 \\
Grantham \(\sim\) ProtSub & +0.54 \\
Miyata \(\sim\) polar requirement & +0.50 \\
Miyata \(\sim\) Kyte--Doolittle & +0.62 \\
Miyata \(\sim\) ProtSub & \ensuremath{-}0.06 \\
Polar requirement \(\sim\) Kyte--Doolittle & +0.01 \\
Polar requirement \(\sim\) ProtSub & +0.08 \\
Kyte--Doolittle \(\sim\) ProtSub & +0.39 \\
\end{longtable}

\protect\phantomsection\label{tbl:s-metric-partial}{}

The picture from \Cref{tbl:s-metric-corr-pairs},
\Cref{tbl:s-metric-corr-codes}, and \Cref{tbl:s-metric-partial} is
consistent with our framing: the metrics share substantial common
variance (median pairwise \(\rho \approx 0.82\) at the code level), but
each contributes unique signal not captured by any combination of the
others. The five-metric convergence is convergent operationalisation,
not redundant replication.

\section{Walsh--Hadamard / 2-adic spectral probe}\label{sec:s-walsh}

This section establishes a finite, computable bridge between our
discrete \({\text{GF}(2)}^{6}\) hypercube framework and the 2-adic codon
algebra developed in the BioSystems literature by
\citep{khrennikov2007}, \citep{dragovich2019}, and
\citep{axelsson2024a}. Reproducible via \texttt{codon-topo} (module
\texttt{src/codon\_topo/analysis/walsh\_2adic.py}; tests
\texttt{tests/test\_walsh\_2adic.py}, 17/17 passing).

For each amino-acid block we form the 0/1 indicator function on the
64-codon space (indexed under the default \(C \rightarrow 00\),
\(U \rightarrow 01\), \(A \rightarrow 10\), \(G \rightarrow 11\)
encoding), compute its Walsh--Hadamard transform (the Fourier transform
on the group \({\text{GF}(2)}^{6}\)), take the 2-adic valuation
\(v_{2}\) of each integer Walsh coefficient, and sum across the spectrum
and across all blocks. We call this aggregate the \textbf{Walsh spectral
depth} of a code.

\subsection{Block-size null}

Under a block-size-matched null (\(n =\)2,000 random partitions with the
same multiset of block sizes as the standard code), the standard code's
spectral depth is anomalously low (\Cref{tbl:s-walsh}). The depth value
is \textbf{mathematically invariant} across all 24 base-to-bit
bijections (the statistic depends only on the underlying partition
geometry, not on the encoding).

\begin{longtable}[]{@{}>{\raggedright\arraybackslash}p{(\linewidth - 8\tabcolsep) * \real{0.45}}>{\raggedleft\arraybackslash}p{(\linewidth - 8\tabcolsep) * \real{0.10}}>{\raggedleft\arraybackslash}p{(\linewidth - 8\tabcolsep) * \real{0.16}}>{\raggedleft\arraybackslash}p{(\linewidth - 8\tabcolsep) * \real{0.10}}>{\raggedleft\arraybackslash}p{(\linewidth - 8\tabcolsep) * \real{0.19}}@{}}
\caption{Walsh spectral depth of the standard code under three null
models. The block-size matched null gives strong shallowness
(\(z = - 17.74\)), but the stricter wobble-box-preserving label
permutation reveals that the depth is mathematically constant under this
null: an algebraic invariant of the (wobble box \(\times\) AA slot
multiset) structure. The Walsh signature of the standard code is fully
encoded in the hierarchical wobble-box decomposition of synonymous
codons.}\tabularnewline
\toprule\noalign{}
\textbf{Null model} & \textbf{Observed} & \textbf{Null mean \(\pm\) SD}
& \textbf{\emph{z}} & \textbf{Frac. \(\leq\) obs.} \\
\midrule\noalign{}
\endfirsthead
\toprule\noalign{}
\textbf{Null model} & \textbf{Observed} & \textbf{Null mean \(\pm\) SD}
& \textbf{\emph{z}} & \textbf{Frac. \(\leq\) obs.} \\
\midrule\noalign{}
\endhead
\bottomrule\noalign{}
\endlastfoot
Block-size matched (\(n =\)2,000) & 544 & 689.3 \(\pm\) 8.2 & \ensuremath{-}17.74 &
0.0000 \\
Wobble-box-preserving label permutation (\(n =\)2,000) & 544 & 544.0
\(\pm\) 0.0 & n/a (invariant) & 1.0000 \\
Encoding sweep (24 encodings \(\times\) 1,500 nulls each) & 544 & n/a
per encoding & {[}\ensuremath{-}18.9, \ensuremath{-}16.9{]} & 0.0000 \\
\end{longtable}

\protect\phantomsection\label{tbl:s-walsh}{}

\begin{figure}
\centering
\emfig{FigW\_walsh\_spectrum}
\caption{2-adic Walsh spectral-depth signature of the standard code (R1
response, engagement with the Khrennikov/Dragovich/Axelsson 2-adic codon
framework). \textbf{Panel A:} block-size-matched null distribution of
the Walsh spectral depth (\(n =\)2,000 random partitions with the same
multiset of block sizes as the standard code). Because per-draw null
samples are not persisted to \texttt{output/walsh\_2adic.json}, the null
is shown as a normal approximation with mean \(= 689.3\) and SD
\(= 8.2\); the standard code's observed depth \(= 544\) (red) sits
\(z = - 17.74\) standard deviations below the null mean (empirical
fraction of null \(\leq\) observed \(= 0\)). \textbf{Panel B:}
encoding-invariance sweep across all \(24\) base-to-bit bijections. Each
encoding is evaluated against its own independent block-size null
(\(n =\)1,500 per encoding, deterministic seed); z-scores fall in the
tight range \(\lbrack - 18.90, - 16.87\rbrack\) (mean \(- 17.87\)), all
\(\approx 5\) times the magnitude of the two-sided \(p = 0.001\)
threshold (\(|z| = 3.29\), dashed red line; note that the
\(z\)-magnitude ratio is not a significance comparison and is quoted
only as a descriptive scale). The observed spectral depth itself equals
\(544\) under every encoding (mathematically invariant; see §S21.3).
Together the two panels quantify the descriptive bridge to the 2-adic
literature that R1 requested: the standard code's Walsh signature is
anomalously shallow versus a block-size-matched null and is a property
of the code, not of any particular base-to-bit bijection.}
\end{figure}

\protect\phantomsection\label{fig:s-walsh-spectrum}{}

\subsection{Wobble-box-preserving label-permutation null}

The block-size-matched null does not control for the first-two-base
wobble box structure that synonymous codons share. We additionally
constructed a stricter null that fixes the 16 first-two-base boxes and
each box's internal partition shape, randomising only the AA labels
assigned to box-slots and preserving each AA's multiset of slot sizes
from the standard code. Under this null the spectral depth statistic is
mathematically constant at 544 (\Cref{tbl:s-walsh}, row 2). This is
consistent with the cautionary remark anticipated in the module's
docstring (\texttt{WOBBLE\_CAVEAT}): the spectral-depth statistic is
invariant under several ``natural'' randomisations of the wobble
structure, and the present null is one of them.

We do \textbf{not} interpret this as evidence of beyond-wobble
optimisation. The honest reading is that the standard code's 2-adic
Walsh signature is an algebraic invariant of the wobble-box
decomposition, which is itself the hierarchical structure that
\citep{khrennikov2007} and follow-up work in the BioSystems literature
model continuously on \(\mathbb{Z}_{2}\). The Walsh--Hadamard / 2-adic
probe is therefore a methodological bridge between the discrete
(hypercube) and continuous (p-adic) descriptions of codon space, not a
separate optimality test.

\subsection{Encoding invariance}

Across all 24 base-to-bit bijections, the standard code's spectral depth
equals 544 (mathematically invariant). The \(z\)-score against the
block-size matched null also remains in the range
\(\lbrack - 18.9, - 16.9\rbrack\), mean \(- 17.9\), across all encodings
(each evaluated with its own independent null draw from a deterministic
seed). \Cref{fig:s-walsh-spectrum} visualises both the block-size null
density (panel A) and the encoding-invariance sweep (panel B).

\subsection{A second Walsh invariant: the wobble-free label spectrum}

The block-indicator spectral depth (§S21.1--§S21.2) characterises the
\textbf{geometry} of the synonymous partition. A complementary Walsh
statistic characterises the \textbf{amino-acid labeling} attached to
that geometry. For each sense amino acid \(a\) let
\(f_{a} \in \left\{ 0,1 \right\}^{64}\) be its indicator vector and
\(F_{a}\) its Walsh--Hadamard transform; partition the \(63\) non-DC
Walsh frequencies \(y \in \lbrack 1,64)\) into a \textbf{wobble-free}
layer (frequencies with zero support on the two position-3 wobble bits;
\(15\) frequencies) and a \textbf{wobble-active} layer (the remaining
\(48\)). The label-spectrum fraction is

\[S = \frac{\sum_{a \in \mathcal{A}\backslash\left\{ \text{Stop} \right\}}\sum_{y\text{ wobble-free},y \neq 0}\left| {F_{a}(y)} \right|^{2}}{\sum_{a \in \mathcal{A}\backslash\left\{ \text{Stop} \right\}}\sum_{y \neq 0}\left| {F_{a}(y)} \right|^{2}}\]
\protect\phantomsection\label{eq:S_label}{}

For the standard code, \(S = 0.7514309076\), mathematically invariant
across all \(24\) base-to-bit bijections
(\texttt{output/walsh\_label\_spectrum.json}).

\textbf{Two nulls, two distinct findings.} Under the \textbf{free}
label-permutation null preserving only the amino-acid count multiset
(stops fixed; \(n =\)10,000, seed \(135325\)): null mean \(= 0.2383\),
sd \(= 0.0169\), range \(\lbrack 0.2020,0.3328\rbrack\),
\(z = + 30.33\), \(0\) of 10,000 draws reach the observed value
(empirical \(p < \frac{1}{n + 1} \approx 10^{- 4}\)). Under the stricter
\textbf{wobble-box-preserving} label-permutation null (the same null
used for \texttt{spectral\_depth} in §S21.2: fixes the \(16\)
first-two-base boxes and each box's internal partition shape, randomises
only which amino-acid label is assigned to each box-slot, preserves each
AA's slot multiset): \(S = 0.7514309076\) in every draw --
mathematically invariant.

\textbf{Honest interpretation.} The free-permutation null shatters the
wobble-box structure of \(4\)-fold-degenerate amino acids, scattering
their Walsh energy into wobble-active frequencies; the box-preserving
null does not. The extreme \(z = + 30\) against the free null therefore
reflects \textbf{exactly} the wobble-box alignment property of the
standard code, not a separate optimisation layer. We do not claim
``beyond wobble'' optimisation. Rather, \(S\) provides an exact
algebraic signature of the wobble phenomenon in the Walsh basis:
biological wobble degeneracy maps directly to spectral concentration on
the wobble-free layer. The free-null calibration measures spectral
concentration against a structurally naive baseline; it is not a test of
selection beyond wobble.

\textbf{Coset quantisation of the null distribution.} The free
label-permutation null is highly quantised: in 10,000 draws it produced
only \(15\) distinct values of \(S\). This reflects the coarse
interaction between the AA degeneracy classes (\(1\), \(2\), \(3\),
\(4\), \(6\)) and the wobble-coset partition of the Walsh frequency
domain: \(S\) depends only on how those classes intersect the wobble
cosets, not on the within-class label assignment. The observed
\(S = 0.7514\) lay above all sampled null values, so the empirical
upper-tail probability is bounded by the Monte Carlo resolution
(\(p \leq 1/(n + 1) \approx 10^{- 4}\)).

\textbf{Mechanistic decomposition.} Each \(4\)-fold-degenerate amino
acid that occupies a complete wobble box (Pro, Ala, Thr, Val, Gly, plus
the CUx, CGx, UCx cores of Leu, Arg, Ser) contributes \(\approx 100\)\%
of its Walsh energy to the wobble-free layer: its indicator function is
constant on a wobble coset, so its Walsh transform is supported
exclusively on frequencies with zero support on the two position-3
wobble bits. Singleton amino acids (Met, Trp) contribute the baseline
\(15/63 \approx 0.238\); a single \(\delta\)-function has uniformly
distributed Walsh energy. The intermediate \((2,2)\)-split and
\((4,2)\)-split blocks contribute mixed weights. The aggregate
\(S = 0.7514\) is therefore in effect a ``weighted box-faithfulness of
the labeling.''

\textbf{Final framing for the bridge to the \(2\)-adic literature.} The
standard genetic code has two encoding-invariant Walsh signatures of
wobble structure: spectral depth \(= 544\) captures the synonymous-block
geometry (§S21.1), and \(S = 0.7514\) captures the amino-acid labeling's
concentration on wobble-free frequency layers. Both support the
descriptive bridge to \(2\)-adic descriptions of the genetic code
\citep{khrennikov2007, axelsson2024a} neither is evidence of selection
beyond wobble.

\section{Slavov (Tsour et al. 2026) SAAP
cross-analysis}\label{sec:s-slavov}

\citep{tsour2026} published an LC-MS proteomic survey identifying 8,746
unique amino-acid substitutions across 1,767 genes in human and mouse
tissues arising from alternate RNA decoding. Their high-confidence
subset (Supplementary Data 8 of \citep{tsour2026} positional probability
\(> 0.9\)) contains 5,873 events spanning 178 unique amino-acid
substitution types. We tested whether the observed substitution events
are enriched for low-codon-Hamming-distance pairs, as predicted by the
codon--anticodon mismatch mechanism that underlies our
\({\text{GF}(2)}^{6}\) adjacency framework.

For each event we computed the minimum nucleotide-Hamming distance
between any source codon (in the standard code) and any target codon.
Filtering to unambiguous 20-AA pairs left 5,611 events and 166 unique
observed substitution types. \Cref{tbl:s-slavov} tabulates the
distance-1/2/3 event counts and baseline fractions;
\Cref{fig:s-slavov-saap} plots the same distribution as grouped bars
for visual comparison.

\begin{longtable}[]{@{}lrrr@{}}
\caption{Distribution of minimum nucleotide-Hamming distance between
source and target codons for the 5,611 high-confidence
alternate-translation events of \citep{tsour2026} (Supplementary Data
8). Single-nucleotide-distance substitutions are dramatically
over-represented relative to the all-380-pair baseline (binomial test,
\(p < 10^{- 100}\)). This is the empirical signature of codon--anticodon
mismatch decoding (Hamming-1) on top of lower-rate mechanisms (RNA
modifications, near-cognate decoding) that produce higher-distance
substitutions.}\tabularnewline
\toprule\noalign{}
\textbf{Min codon NT distance} & \textbf{Events} & \textbf{Event \%} &
\textbf{Baseline (380 pairs) \%} \\
\midrule\noalign{}
\endfirsthead
\toprule\noalign{}
\textbf{Min codon NT distance} & \textbf{Events} & \textbf{Event \%} &
\textbf{Baseline (380 pairs) \%} \\
\midrule\noalign{}
\endhead
\bottomrule\noalign{}
\endlastfoot
1 & 3,649 & 65.0\% & 39.5\% \\
2 & 1,062 & 18.9\% & 53.2\% \\
3 & 900 & 16.0\% & 7.4\% \\
\end{longtable}

\protect\phantomsection\label{tbl:s-slavov}{}

\begin{figure}
\centering
\emfig{FigSl\_slavov\_saap}
\caption{Visual companion to \Cref{tbl:s-slavov} (added in response to
R3, who noted that a three-row table understates the effect). Grouped
bars compare the observed nucleotide-Hamming-distance distribution of
Tsour et al.'s 5,611 high-confidence sense-codon recoding events (blue)
against the null distribution across all 380 amino-acid-to-amino-acid
pairs in the standard code (grey). At distance 1 the observed events are
65.0\% vs a 39.5\% baseline (event-weighted binomial \(p < 10^{- 100}\);
\citep{tsour2026}); at distance 2 the observed events are
\textbf{depleted} (18.9\% vs 53.2\%); at distance 3 they are moderately
enriched (16.0\% vs 7.4\%). The Hamming-1 concentration is the empirical
signature of codon--anticodon mismatch decoding operating on top of
lower-rate mechanisms (RNA modifications, near-cognate decoding) that
produce higher-distance substitutions. Figure and table are rendered
from the same \texttt{output/slavov\_saap\_codon\_distances.json} (event
and baseline distance distributions), so they cannot drift.}
\end{figure}

\protect\phantomsection\label{fig:s-slavov-saap}{}

At the level of substitution \textbf{types} (unique amino-acid pairs
rather than events), the enrichment is modest (41.6\% of observed types
are at min-NT-1 vs 39.5\% baseline; Fisher exact OR = 1.17,
\(p = 0.26\); Mann--Whitney U on NT distances, observed vs unobserved
AAS types: \(p = 0.18\)). This indicates that essentially every
amino-acid substitution type can occur in the Tsour et al. dataset: it
is the \textbf{frequency} of each type that is biased toward
single-nucleotide-distance pairs. The observation matches the
mechanistic inventory reported by \citep{tsour2026}: codon--anticodon
mismatch is the dominant high-frequency mechanism, with RNA
modifications and near-cognate decoding contributing lower-frequency
events at higher codon-Hamming distance. The Hamming-1 enrichment is
independent of which amino-acid distance metric is used and is an
external empirical test of the codon-adjacency assumption underlying the
main-text edge-mismatch objective (equation 1).

\section{Exploratory observations: full detail}\label{sec:s-exploratory}

The main text (§3.7) summarises three exploratory observations without
inferential weight; this section records the underlying detail for each.

\subsection{Bit-position bias in codon
reassignments}\label{sec:s-bitbias}

The distribution of bit-flips across the 6 coordinates of
\({\text{GF}(2)}^{6}\) in natural codon reassignments shows apparent
positional skew under a uniform null (\(\chi^{2} = 16.26\),
\(p = 0.006\), \(\text{df } = 5\)). However, this signal is
substantially attenuated after de-duplication to 20 unique (codon,
target amino acid) pairs (\(p = 0.075\)) and vanishes entirely under a
codon-preserving permutation null (\(p = 1.0\)). The apparent bias is
therefore explained by which codons are recurrently reassigned across
lineages, not by a genuine positional preference in
\({\text{GF}(2)}^{6}\). This is an example of a signal that survives one
null model (uniform) but fails another (codon-preserving), and is
retained in the record only to document the diagnostic sequence.

\begin{figure}
\centering
\emfig{FigG\_bit\_position\_bias}
\caption{Bit-position bias in codon-reassignment events. Observed
bit-flip counts at each of the 6 coordinates of \({\text{GF}(2)}^{6}\),
overlaid on the uniform-null expectation and the codon-preserving
permutation null. The apparent uniform-null skew (\(p = 0.006\)) is
fully explained by the codon-preserving null (\(p = 1.0\)), consistent
with a recurrent-codon effect rather than a genuine positional
preference. Retained as a null-model-diagnostic example.}
\end{figure}

\protect\phantomsection\label{fig:s-bit-position-bias}{}

\subsection{Variant-code disconnection catalogue}\label{sec:s-catalogue}

A systematic survey across all 27 NCBI translation tables identifies
four lineage-collapsed variant-code amino-acid disconnections at
\(\varepsilon = 1\) in \({\text{GF}(2)}^{6}\) under the default
encoding: threonine in the yeast mitochondrial code (translation table
3, CUN\(\rightarrow\)Thr); leucine in the chlorophycean mitochondrial
codes (translation tables 16 and 22, both with UAG\(\rightarrow\)Leu;
table 16 is the chlorophycean mitochondrial code \citep{hayashi1996} and
table 22 is the closely related \emph{Scenedesmus obliquus}
mitochondrial code, which additionally reassigns UCA
Ser\(\rightarrow\)Stop, and both produce equivalent \(\varepsilon = 2\)
Leu reconnection profiles so they collapse to a single
algal-mitochondrial event); alanine in \emph{Pachysolen tannophilus}
nuclear code (translation table 26, CUG\(\rightarrow\)Ala); and a
tripartite serine in the \emph{Candida}-clade alternative yeast nuclear
code (translation table 12, CUG\(\rightarrow\)Ser; \citep{santos1999}).
These cases, combined with the universal serine disconnection, make up
the complete inventory of amino-acid graph disconnections at unit
Hamming distance under the default encoding.

A separate and weaker geometric exception (specific to filtration rather
than to disconnection) arises in translation table 32 (Balanophoraceae
plastid; UAG\(\rightarrow\)Trp). Trp is 1-fold (UGG only) under the
standard code; in table 32 it becomes 2-fold (UGG, UAG). The pair
differs in the second nucleotide (G\(\leftrightarrow\)A), i.e. at bit
position 3 in our 0-based 6-bit indexing (the second bit of the second
nucleotide), not at the wobble bit (position 5) where every
standard-code 2-fold pair sits. Hamming distance 1, so the pair is
connected at \(\varepsilon = 1\) and adds no new entry to the
disconnection catalogue. An empirical scan over every 2-fold amino-acid
pair in all 27 NCBI tables confirms that this is the unique deviation
from the bit-5 two-fold filtration: the standard code's nine 2-fold
amino acids and every analogous pair introduced by the other 26 tables
differ at bit 5.

\subsection{Atchley Factor 3 and Serine
convergence}\label{sec:s-atchley}

Serine has the most extreme Atchley Factor 3 score among the 20 amino
acids (\(F_{3} = - 4.760\), 2.24 SD below the mean;
\citep{atchley2005}), and it is the only amino acid disconnected at
\(\varepsilon = 1\) under every base-to-bit encoding. This convergence
is unsurprising: Atchley \(F_{3}\) is a composite of molecular size and
codon diversity that partly captures codon-family structure, so the
\({\text{GF}(2)}^{6}\) topology and \(F_{3}\) are not fully independent
views of Serine's anomaly. Both reflect Serine's disproportionate codon
diversity (6 codons in two disconnected families) relative to its small
physicochemical footprint; the \({\text{GF}(2)}^{6}\) framework provides
a complementary structural account rather than independent
corroboration.

\subsection{Three-tier tRNA mechanistic landscape (from main-text §3.6)}

The manuscript §3.6 summarises the tRNA enrichment result; the
mechanistic detail moved from that section is recorded here for
completeness. The largest observed case is \emph{Tetrahymena
thermophila} (NCBI translation table 6, UAA/UAG reassigned to Gln;
\citep{hanyu1986}), which carries 54 glutamine tRNA genes (including 39
suppressor tRNAs reading the reassigned stop codons), compared to 3 Gln
tRNAs in the standard-code ciliate \emph{Ichthyophthirius multifiliis}
(assembly GCF\_000220395.1), 11 in \emph{Stentor coeruleus} (assembly
GCA\_001970955.1), and 3 in \emph{Fabrea salina} (assembly
GCA\_022984795.1 from \citep{zhang2022}). All four counts were generated
in this work by running tRNAscan-SE 2.0.12 on the listed assemblies. The
pattern extends across reassignment types. Among Gln-reassignment
ciliates, \emph{Pseudocohnilembus persalinus} (20 Gln tRNAs including 15
suppressors) and \emph{Halteria grandinella} (9 Gln tRNAs including 3
suppressors; \citep{zheng2021}) represent independent lineages within
Oligohymenophorea and Spirotrichea respectively. Among Cys-reassignment
ciliates (NCBI translation table 10, UGA\(\rightarrow\)Cys), six
tRNAscan-SE--verified \emph{Euplotes} species all carry tRNA-Cys genes
with the non-canonical TCA anticodon (reading UGA), with 1--4 such genes
per species alongside standard GCA-anticodon Cys tRNAs.

However, the pattern is not universal. \emph{Blastocrithidia nonstop}
(NCBI translation table 31) reassigned all three stop codons (as did
\emph{Condylostoma magnum}, where stop-codon function is
context-dependent; \citep{heaphy2016}) but achieved
UGA\(\rightarrow\)Trp via anticodon stem shortening (5 bp
\(\rightarrow\) 4 bp) of tRNA-Trp(CCA), combined with an eRF1 Ser74Gly
mutation, rather than gene duplication \citep{kachale2023}. Similarly,
\emph{Mycoplasmoides} species with UGA\(\rightarrow\)Trp use a single
tRNA-Trp with anticodon modification. These boundary cases define a
three-tier mechanistic landscape: (i) tRNA gene duplication in large
nuclear genomes, (ii) anticodon structural modification in streamlined
genomes, and (iii) anticodon base modification in minimal genomes.

\section{Software and reproducibility}\label{sec:s-software}

This section provides the metadata needed to reproduce every number in
the manuscript and supplement \textbf{within numerical and rendering
tolerance} from the public repository. Every figure, table, and inline
statistic is rendered by the Typst sources \texttt{manuscript.typ} and
\texttt{supplement.typ} (also in the repository) from the JSON outputs
of a single \texttt{codon-topo\ all} invocation, so the manuscript and
supplement cannot drift from each other within a single pipeline run. We
do not claim bit-for-bit reproducibility: the pinned software floor
below fixes randomness (seeded RNG) and language versions, but does not
pin every transitive dependency, every optimizer default, or every
renderer version, so exact byte-level identity of PDFs, PNG rasters, and
floating-point outputs across systems is not guaranteed. Numerical
values in the JSON artifacts, quantile ranks, and permutation-test
\(p\)-values should reproduce to at least three significant figures;
rendered PDFs may differ in font subsetting, image compression, and
Typst layout across Typst versions.

All analyses were performed using the \texttt{codon-topo} Python package
(version 0.6.1, tag \texttt{v0.6.1}). The code is publicly released at
\url{https://github.com/biostochastics/codontopo.} Dependencies and
runtime requirements:

\begin{itemize}
\item
  Python 3.11 (tested on 3.11.14), NumPy 1.24+, SciPy 1.10+
\item
  R 4.5, ggplot2, ggpubr, viridis, patchwork (for figures)
\item
  tRNAscan-SE 2.0.12 with Infernal 1.1.4
\item
  Typst 0.14.2 (for manuscript and supplement PDFs)
\item
  Random seed: 135325 (all Monte Carlo analyses)
\item
  Exact package versions used in the reference build are listed in
  \texttt{requirements.lock} (Python) and \texttt{renv.lock} (R) in the
  release tag; a \texttt{Dockerfile} is provided for a containerized
  rebuild. Users who install unpinned dependencies should expect
  reproducibility only within numerical/rendering tolerance rather than
  bit-for-bit.
\end{itemize}

Analyses are fully reproducible via:

\begin{verbatim}
git clone https://github.com/biostochastics/codontopo.git
cd codontopo
pip install -e ".[all]"
codon-topo all --output-dir=./output --seed=135325
python scripts/finalize_manuscript_stats.py
python scripts/generate_tables.py
Rscript src/codon_topo/visualization/R/strengthened_figures.R
\end{verbatim}

The \texttt{{[}all{]}} extra installs every dependency required to
reproduce the analyses; the \texttt{{[}dev{]}} extra also installs the
test toolchain. The complete test suite (432 tests) is then runnable
with \texttt{python3.11\ -m\ pytest\ tests/\ -m\ "not\ slow"} after
\texttt{pip\ install\ -e\ ".{[}dev{]}"}.


% Corrected 2026-08-15: bibliography inlined directly (see
% manuscript.tex for rationale).
\begin{thebibliography}{25}
\expandafter\ifx\csname natexlab\endcsname\relax\def\natexlab#1{#1}\fi
\providecommand{\url}[1]{\texttt{#1}}
\providecommand{\href}[2]{#2}
\providecommand{\path}[1]{#1}
\providecommand{\DOIprefix}{doi:}
\providecommand{\ArXivprefix}{arXiv:}
\providecommand{\URLprefix}{URL: }
\providecommand{\Pubmedprefix}{pmid:}
\providecommand{\doi}[1]{\href{http://dx.doi.org/#1}{\path{#1}}}
\providecommand{\Pubmed}[1]{\href{pmid:#1}{\path{#1}}}
\providecommand{\bibinfo}[2]{#2}
\ifx\xfnm\relax \def\xfnm[#1]{\unskip,\space#1}\fi
%Type = Article
\bibitem[{Antoneli and Forger(2011)}]{antoneli2011}
\bibinfo{author}{Antoneli, F.}, \bibinfo{author}{Forger, M.},
  \bibinfo{year}{2011}.
\newblock \bibinfo{title}{Symmetry breaking in the genetic code: finite
  groups}.
\newblock \bibinfo{journal}{Mathematical and Computer Modelling}
  \bibinfo{volume}{53}, \bibinfo{pages}{1469--1488}.
\newblock \DOIprefix\doi{10.1016/j.mcm.2010.03.050}.
%Type = Article
\bibitem[{Atchley et~al.(2005)Atchley, Zhao, Fernandes and
  Dr{\"u}ke}]{atchley2005}
\bibinfo{author}{Atchley, W.R.}, \bibinfo{author}{Zhao, J.},
  \bibinfo{author}{Fernandes, A.D.}, \bibinfo{author}{Dr{\"u}ke, T.},
  \bibinfo{year}{2005}.
\newblock \bibinfo{title}{Solving the protein sequence metric problem}.
\newblock \bibinfo{journal}{Proceedings of the National Academy of Sciences}
  \bibinfo{volume}{102}, \bibinfo{pages}{6395--6400}.
\newblock \DOIprefix\doi{10.1073/pnas.0408677102}.
%Type = Article
\bibitem[{Buschmann et~al.(2026)Buschmann, Bolz, El-Hendi, Malekian and
  Schroeder}]{buschmann2026}
\bibinfo{author}{Buschmann, L.}, \bibinfo{author}{Bolz, S.N.},
  \bibinfo{author}{El-Hendi, F.}, \bibinfo{author}{Malekian, N.},
  \bibinfo{author}{Schroeder, M.}, \bibinfo{year}{2026}.
\newblock \bibinfo{title}{Much ado about nothing: modeling amino acid
  replacement with predicted protein structures}.
\newblock \bibinfo{journal}{Bioinformatics} \bibinfo{volume}{42},
  \bibinfo{pages}{btag188}.
\newblock \DOIprefix\doi{10.1093/bioinformatics/btag188}.
%Type = Article
\bibitem[{Chan and Lowe(2016)}]{chan2016}
\bibinfo{author}{Chan, P.P.}, \bibinfo{author}{Lowe, T.M.},
  \bibinfo{year}{2016}.
\newblock \bibinfo{title}{{GtRNAdb} 2.0: an expanded database of transfer {RNA}
  genes identified in complete and draft genomes}.
\newblock \bibinfo{journal}{Nucleic Acids Research} \bibinfo{volume}{44},
  \bibinfo{pages}{D184--D189}.
\newblock \DOIprefix\doi{10.1093/nar/gkv1309}.
%Type = Article
\bibitem[{Chan and Lowe(2019)}]{chan2019}
\bibinfo{author}{Chan, P.P.}, \bibinfo{author}{Lowe, T.M.},
  \bibinfo{year}{2019}.
\newblock \bibinfo{title}{{tRNAscan-SE}: searching for {tRNA} genes in genomic
  sequences}.
\newblock \bibinfo{journal}{Methods in Molecular Biology}
  \bibinfo{volume}{1962}, \bibinfo{pages}{1--14}.
\newblock \DOIprefix\doi{10.1007/978-1-4939-9173-0_1}.
%Type = Unpublished
\bibitem[{Clayworth(2026)}]{clayworth2026}
\bibinfo{author}{Clayworth, P.}, \bibinfo{year}{2026}.
\newblock \bibinfo{title}{Algebraic structure of the genetic code: a technical
  note}.
\newblock \bibinfo{note}{Companion methodological note. Source code and full
  claim hierarchy are released alongside the present manuscript at
  https://github.com/biostochastics/codontopo; the published manuscript
  supersedes the technical note in scope and statistical detail.}
%Type = Article
\bibitem[{Crick(1966)}]{crick1966}
\bibinfo{author}{Crick, F.H.C.}, \bibinfo{year}{1966}.
\newblock \bibinfo{title}{Codon--anticodon pairing: the wobble hypothesis}.
\newblock \bibinfo{journal}{Journal of Molecular Biology} \bibinfo{volume}{19},
  \bibinfo{pages}{548--555}.
\newblock \DOIprefix\doi{10.1016/S0022-2836(66)80022-0}.
%Type = Article
\bibitem[{Di~Giulio(2001)}]{digiulio2001}
\bibinfo{author}{Di~Giulio, M.}, \bibinfo{year}{2001}.
\newblock \bibinfo{title}{The origin of the genetic code cannot be studied
  using measurements based on the {PAM} matrix because this matrix reflects the
  code itself, making any such analyses tautologous}.
\newblock \bibinfo{journal}{Journal of Theoretical Biology}
  \bibinfo{volume}{208}, \bibinfo{pages}{141--144}.
\newblock \DOIprefix\doi{10.1006/jtbi.2000.2206}.
%Type = Article
\bibitem[{Dragovich and Mi{\v s}i{\'c}(2019)}]{dragovich2019}
\bibinfo{author}{Dragovich, B.}, \bibinfo{author}{Mi{\v s}i{\'c}, N.{\v Z}.},
  \bibinfo{year}{2019}.
\newblock \bibinfo{title}{p-{A}dic hierarchical properties of the genetic
  code}.
\newblock \bibinfo{journal}{BioSystems} \bibinfo{volume}{185},
  \bibinfo{pages}{104017}.
\newblock \DOIprefix\doi{10.1016/j.biosystems.2019.104017}.
%Type = Article
\bibitem[{Freeland and Hurst(1998)}]{freeland1998}
\bibinfo{author}{Freeland, S.J.}, \bibinfo{author}{Hurst, L.D.},
  \bibinfo{year}{1998}.
\newblock \bibinfo{title}{The genetic code is one in a million}.
\newblock \bibinfo{journal}{Journal of Molecular Evolution}
  \bibinfo{volume}{47}, \bibinfo{pages}{238--248}.
\newblock \DOIprefix\doi{10.1007/PL00006381}.
%Type = Article
\bibitem[{Grantham(1974)}]{grantham1974}
\bibinfo{author}{Grantham, R.}, \bibinfo{year}{1974}.
\newblock \bibinfo{title}{Amino acid difference formula to help explain protein
  evolution}.
\newblock \bibinfo{journal}{Science} \bibinfo{volume}{185},
  \bibinfo{pages}{862--864}.
\newblock \DOIprefix\doi{10.1126/science.185.4154.862}.
%Type = Article
\bibitem[{Haig and Hurst(1991)}]{haig1991}
\bibinfo{author}{Haig, D.}, \bibinfo{author}{Hurst, L.D.},
  \bibinfo{year}{1991}.
\newblock \bibinfo{title}{A quantitative measure of error minimization in the
  genetic code}.
\newblock \bibinfo{journal}{Journal of Molecular Evolution}
  \bibinfo{volume}{33}, \bibinfo{pages}{412--417}.
\newblock \DOIprefix\doi{10.1007/BF02103132}.
%Type = Article
\bibitem[{Hanyu et~al.(1986)Hanyu, Kuchino, Nishimura and Beier}]{hanyu1986}
\bibinfo{author}{Hanyu, N.}, \bibinfo{author}{Kuchino, Y.},
  \bibinfo{author}{Nishimura, S.}, \bibinfo{author}{Beier, H.},
  \bibinfo{year}{1986}.
\newblock \bibinfo{title}{Dramatic events in ciliate evolution: alteration of
  {UAA} and {UAG} termination codons to glutamine codons due to anticodon
  mutations in two {Tetrahymena} trnas-gln}.
\newblock \bibinfo{journal}{The EMBO Journal} \bibinfo{volume}{5},
  \bibinfo{pages}{1307--1311}.
\newblock \DOIprefix\doi{10.1002/j.1460-2075.1986.tb04360.x}.
%Type = Article
\bibitem[{Hayashi-Ishimaru et~al.(1996)Hayashi-Ishimaru, Ohama, Kawatsu,
  Nakamura and Osawa}]{hayashi1996}
\bibinfo{author}{Hayashi-Ishimaru, Y.}, \bibinfo{author}{Ohama, T.},
  \bibinfo{author}{Kawatsu, Y.}, \bibinfo{author}{Nakamura, K.},
  \bibinfo{author}{Osawa, S.}, \bibinfo{year}{1996}.
\newblock \bibinfo{title}{{UAG} is a sense codon in several chlorophycean
  mitochondria}.
\newblock \bibinfo{journal}{Current Genetics} \bibinfo{volume}{30},
  \bibinfo{pages}{29--33}.
\newblock \DOIprefix\doi{10.1007/s002940050096}. \bibinfo{note}{nCBI
  translation table 16 (chlorophycean mitochondrial); UAG $\to$ Leu
  reassignment in green-algal mitochondria}.
%Type = Article
\bibitem[{Heaphy et~al.(2016)Heaphy, Mariotti, Gladyshev, Atkins and
  Baranov}]{heaphy2016}
\bibinfo{author}{Heaphy, S.M.}, \bibinfo{author}{Mariotti, M.},
  \bibinfo{author}{Gladyshev, V.N.}, \bibinfo{author}{Atkins, J.F.},
  \bibinfo{author}{Baranov, P.V.}, \bibinfo{year}{2016}.
\newblock \bibinfo{title}{Novel ciliate genetic code variants including the
  reassignment of all three stop codons to sense codons in {Condylostoma
  magnum}}.
\newblock \bibinfo{journal}{Molecular Biology and Evolution}
  \bibinfo{volume}{33}, \bibinfo{pages}{2885--2889}.
\newblock \DOIprefix\doi{10.1093/molbev/msw166}.
%Type = Article
\bibitem[{Kachale et~al.(2023)Kachale, Pavl{\'\i}kov{\'a}, Nenarokova,
  Roithov{\'a}, Durante, Milet{\'\i}nov{\'a}, Z{\'a}honov{\'a}, Nenarokov,
  Vot{\'y}pka, Hor{\'a}kov{\'a}, Ross, Yurchenko, Beznoskov{\'a}, Paris,
  Val{\'a}{\v s}ek and Luke{\v s}}]{kachale2023}
\bibinfo{author}{Kachale, A.}, \bibinfo{author}{Pavl{\'\i}kov{\'a}, Z.},
  \bibinfo{author}{Nenarokova, A.}, \bibinfo{author}{Roithov{\'a}, A.},
  \bibinfo{author}{Durante, I.M.}, \bibinfo{author}{Milet{\'\i}nov{\'a}, P.},
  \bibinfo{author}{Z{\'a}honov{\'a}, K.}, \bibinfo{author}{Nenarokov, S.},
  \bibinfo{author}{Vot{\'y}pka, J.}, \bibinfo{author}{Hor{\'a}kov{\'a}, E.},
  \bibinfo{author}{Ross, R.L.}, \bibinfo{author}{Yurchenko, V.},
  \bibinfo{author}{Beznoskov{\'a}, P.}, \bibinfo{author}{Paris, Z.},
  \bibinfo{author}{Val{\'a}{\v s}ek, L.S.}, \bibinfo{author}{Luke{\v s}, J.},
  \bibinfo{year}{2023}.
\newblock \bibinfo{title}{Short {tRNA} anticodon stem and mutant {eRF1} allow
  stop codon reassignment}.
\newblock \bibinfo{journal}{Nature} \bibinfo{volume}{613},
  \bibinfo{pages}{751--758}.
\newblock \DOIprefix\doi{10.1038/s41586-022-05584-2}.
%Type = Article
\bibitem[{Khrennikov and Kozyrev(2007)}]{khrennikov2007}
\bibinfo{author}{Khrennikov, A.Y.}, \bibinfo{author}{Kozyrev, S.V.},
  \bibinfo{year}{2007}.
\newblock \bibinfo{title}{Genetic code on the diadic plane}.
\newblock \bibinfo{journal}{Physica A: Statistical Mechanics and its
  Applications} \bibinfo{volume}{381}, \bibinfo{pages}{265--272}.
\newblock \DOIprefix\doi{10.1016/j.physa.2007.03.018}.
%Type = Article
\bibitem[{Santos et~al.(1999)Santos, Cheesman, Costa, Moradas-Ferreira and
  Tuite}]{santos1999}
\bibinfo{author}{Santos, M.A.S.}, \bibinfo{author}{Cheesman, C.},
  \bibinfo{author}{Costa, V.}, \bibinfo{author}{Moradas-Ferreira, P.},
  \bibinfo{author}{Tuite, M.F.}, \bibinfo{year}{1999}.
\newblock \bibinfo{title}{Selective advantages created by codon ambiguity
  allowed for the evolution of an alternative genetic code in {Candida} spp.}
\newblock \bibinfo{journal}{Molecular Microbiology} \bibinfo{volume}{31},
  \bibinfo{pages}{937--947}.
\newblock \DOIprefix\doi{10.1046/j.1365-2958.1999.01233.x}.
%Type = Article
\bibitem[{Sengupta et~al.(2007)Sengupta, Yang and Higgs}]{sengupta2007}
\bibinfo{author}{Sengupta, S.}, \bibinfo{author}{Yang, X.},
  \bibinfo{author}{Higgs, P.G.}, \bibinfo{year}{2007}.
\newblock \bibinfo{title}{The mechanisms of codon reassignments in
  mitochondrial genetic codes}.
\newblock \bibinfo{journal}{Journal of Molecular Evolution}
  \bibinfo{volume}{64}, \bibinfo{pages}{662--688}.
\newblock \DOIprefix\doi{10.1007/s00239-006-0284-7}.
%Type = Article
\bibitem[{Singh et~al.(2023)Singh, Seah, Emmerich, Singh, Woehle, Huettel,
  Byerly, Stover, Sugiura, Harumoto and Swart}]{singh2023}
\bibinfo{author}{Singh, M.}, \bibinfo{author}{Seah, B.K.B.},
  \bibinfo{author}{Emmerich, C.}, \bibinfo{author}{Singh, A.},
  \bibinfo{author}{Woehle, C.}, \bibinfo{author}{Huettel, B.},
  \bibinfo{author}{Byerly, A.}, \bibinfo{author}{Stover, N.A.},
  \bibinfo{author}{Sugiura, M.}, \bibinfo{author}{Harumoto, T.},
  \bibinfo{author}{Swart, E.C.}, \bibinfo{year}{2023}.
\newblock \bibinfo{title}{Origins of genome-editing excisases as illuminated by
  the somatic genome of the ciliate {Blepharisma}}.
\newblock \bibinfo{journal}{Proceedings of the National Academy of Sciences}
  \bibinfo{volume}{120}, \bibinfo{pages}{e2213887120}.
\newblock \DOIprefix\doi{10.1073/pnas.2213887120}. \bibinfo{note}{reports the
  \emph{B.~stoltei} ATCC 30299 MAC genome assembly (GCA\_965603825.1; ENA
  Bstoltei\_ATCC30299\_MAC\_1.00) and documents that this strain uses an
  alternative nuclear genetic code with UGA codons reassigned from stops to
  tryptophan (main text and SI Appendix Fig.~S2B), in addition to the canonical
  NCBI table-15 UAG $\rightarrow$ Gln assignment.}
%Type = Article
\bibitem[{Tsour et~al.(2026)Tsour, Machn{\'e}, Leduc, Widmer, Koo, Guez,
  Karczewski and Slavov}]{tsour2026}
\bibinfo{author}{Tsour, S.}, \bibinfo{author}{Machn{\'e}, R.},
  \bibinfo{author}{Leduc, A.}, \bibinfo{author}{Widmer, S.},
  \bibinfo{author}{Koo, E.}, \bibinfo{author}{Guez, J.},
  \bibinfo{author}{Karczewski, K.J.}, \bibinfo{author}{Slavov, N.},
  \bibinfo{year}{2026}.
\newblock \bibinfo{title}{Alternate {RNA} decoding results in stable and
  abundant proteins in mammals}.
\newblock \bibinfo{journal}{Nature} \DOIprefix\doi{10.1038/s41586-026-10678-2}.
%Type = Article
\bibitem[{Yurova~Axelsson and Khrennikov(2024)}]{axelsson2024a}
\bibinfo{author}{Yurova~Axelsson, E.}, \bibinfo{author}{Khrennikov, A.Y.},
  \bibinfo{year}{2024}.
\newblock \bibinfo{title}{Generation of genetic codes with 2-adic codon algebra
  and adaptive dynamics}.
\newblock \bibinfo{journal}{BioSystems} \bibinfo{volume}{240},
  \bibinfo{pages}{105230}.
\newblock \DOIprefix\doi{10.1016/j.biosystems.2024.105230}.
%Type = Article
\bibitem[{Zehir et~al.(2017)Zehir, Benayed, Shah, Syed, Middha, Kim,
  Srinivasan, Gao, Chakravarty, Devlin, Hellmann, Barron, Schram, Hameed,
  Dogan, Ross, Hechtman, DeLair, Yao, Mandelker, Cheng, Chandramohan, Mohanty,
  Ptashkin, Jayakumaran, Prasad, Syed, Rema, Liu, Nafa, Borsu, Sadowska,
  Casanova, Bacares, Kiecka, Razumova, Son, Stewart, Baldi, Mullaney,
  Al-Ahmadie, Vakiani, Abeshouse, Penson, Jonsson, Camacho, Chang, Won, Gross,
  Kundra, Heins, Chen, Phillips, Zhang, Wang, Ochoa, Wills, Eubank, Thomas,
  Gardos, Reales, Galle, Durany, Cambria, Abida, Cercek, Feldman, Gounder,
  Hakimi, Harding, Iyer, Janjigian, Jordan, Kelly, Lowery, Morris, Omuro, Raj,
  Razavi, Shoushtari, Shukla, Soumerai, Varghese, Yaeger, Coleman, Bochner,
  Riely, Saltz, Scher, Sabbatini, Robson, Klimstra, Taylor, Baselga, Schultz,
  Hyman, Arcila, Solit, Ladanyi and Berger}]{zehir2017}
\bibinfo{author}{Zehir, A.}, \bibinfo{author}{Benayed, R.},
  \bibinfo{author}{Shah, R.H.}, \bibinfo{author}{Syed, A.},
  \bibinfo{author}{Middha, S.}, \bibinfo{author}{Kim, H.R.},
  \bibinfo{author}{Srinivasan, P.}, \bibinfo{author}{Gao, J.},
  \bibinfo{author}{Chakravarty, D.}, \bibinfo{author}{Devlin, S.M.},
  \bibinfo{author}{Hellmann, M.D.}, \bibinfo{author}{Barron, D.A.},
  \bibinfo{author}{Schram, A.M.}, \bibinfo{author}{Hameed, M.},
  \bibinfo{author}{Dogan, S.}, \bibinfo{author}{Ross, D.S.},
  \bibinfo{author}{Hechtman, J.F.}, \bibinfo{author}{DeLair, D.F.},
  \bibinfo{author}{Yao, J.}, \bibinfo{author}{Mandelker, D.L.},
  \bibinfo{author}{Cheng, D.T.}, \bibinfo{author}{Chandramohan, R.},
  \bibinfo{author}{Mohanty, A.S.}, \bibinfo{author}{Ptashkin, R.N.},
  \bibinfo{author}{Jayakumaran, G.}, \bibinfo{author}{Prasad, M.},
  \bibinfo{author}{Syed, M.H.}, \bibinfo{author}{Rema, A.B.},
  \bibinfo{author}{Liu, Z.Y.}, \bibinfo{author}{Nafa, K.},
  \bibinfo{author}{Borsu, L.}, \bibinfo{author}{Sadowska, J.},
  \bibinfo{author}{Casanova, J.}, \bibinfo{author}{Bacares, R.},
  \bibinfo{author}{Kiecka, I.J.}, \bibinfo{author}{Razumova, A.},
  \bibinfo{author}{Son, J.B.}, \bibinfo{author}{Stewart, L.},
  \bibinfo{author}{Baldi, T.}, \bibinfo{author}{Mullaney, K.A.},
  \bibinfo{author}{Al-Ahmadie, H.}, \bibinfo{author}{Vakiani, E.},
  \bibinfo{author}{Abeshouse, A.A.}, \bibinfo{author}{Penson, A.V.},
  \bibinfo{author}{Jonsson, P.}, \bibinfo{author}{Camacho, N.},
  \bibinfo{author}{Chang, M.T.}, \bibinfo{author}{Won, H.H.},
  \bibinfo{author}{Gross, B.E.}, \bibinfo{author}{Kundra, R.},
  \bibinfo{author}{Heins, Z.J.}, \bibinfo{author}{Chen, H.W.},
  \bibinfo{author}{Phillips, S.}, \bibinfo{author}{Zhang, H.},
  \bibinfo{author}{Wang, J.}, \bibinfo{author}{Ochoa, A.},
  \bibinfo{author}{Wills, J.}, \bibinfo{author}{Eubank, M.},
  \bibinfo{author}{Thomas, S.B.}, \bibinfo{author}{Gardos, S.M.},
  \bibinfo{author}{Reales, D.N.}, \bibinfo{author}{Galle, J.},
  \bibinfo{author}{Durany, R.}, \bibinfo{author}{Cambria, R.},
  \bibinfo{author}{Abida, W.}, \bibinfo{author}{Cercek, A.},
  \bibinfo{author}{Feldman, D.R.}, \bibinfo{author}{Gounder, M.M.},
  \bibinfo{author}{Hakimi, A.A.}, \bibinfo{author}{Harding, J.J.},
  \bibinfo{author}{Iyer, G.}, \bibinfo{author}{Janjigian, Y.Y.},
  \bibinfo{author}{Jordan, E.J.}, \bibinfo{author}{Kelly, C.M.},
  \bibinfo{author}{Lowery, M.A.}, \bibinfo{author}{Morris, L.G.T.},
  \bibinfo{author}{Omuro, A.M.}, \bibinfo{author}{Raj, N.},
  \bibinfo{author}{Razavi, P.}, \bibinfo{author}{Shoushtari, A.N.},
  \bibinfo{author}{Shukla, N.}, \bibinfo{author}{Soumerai, T.E.},
  \bibinfo{author}{Varghese, A.M.}, \bibinfo{author}{Yaeger, R.},
  \bibinfo{author}{Coleman, J.}, \bibinfo{author}{Bochner, B.},
  \bibinfo{author}{Riely, G.J.}, \bibinfo{author}{Saltz, L.B.},
  \bibinfo{author}{Scher, H.I.}, \bibinfo{author}{Sabbatini, P.J.},
  \bibinfo{author}{Robson, M.E.}, \bibinfo{author}{Klimstra, D.S.},
  \bibinfo{author}{Taylor, B.S.}, \bibinfo{author}{Baselga, J.},
  \bibinfo{author}{Schultz, N.}, \bibinfo{author}{Hyman, D.M.},
  \bibinfo{author}{Arcila, M.E.}, \bibinfo{author}{Solit, D.B.},
  \bibinfo{author}{Ladanyi, M.}, \bibinfo{author}{Berger, M.F.},
  \bibinfo{year}{2017}.
\newblock \bibinfo{title}{Mutational landscape of metastatic cancer revealed
  from prospective clinical sequencing of 10,000 patients}.
\newblock \bibinfo{journal}{Nature Medicine} \bibinfo{volume}{23},
  \bibinfo{pages}{703--713}.
\newblock \DOIprefix\doi{10.1038/nm.4333}.
%Type = Article
\bibitem[{Zhang et~al.(2022)Zhang, Hou, Qi, Hou, Zhang, Zhao and
  Miao}]{zhang2022}
\bibinfo{author}{Zhang, B.}, \bibinfo{author}{Hou, L.}, \bibinfo{author}{Qi,
  H.}, \bibinfo{author}{Hou, L.}, \bibinfo{author}{Zhang, T.},
  \bibinfo{author}{Zhao, F.}, \bibinfo{author}{Miao, M.}, \bibinfo{year}{2022}.
\newblock \bibinfo{title}{An extremely streamlined macronuclear genome in the
  free-living protozoan {Fabrea salina}}.
\newblock \bibinfo{journal}{Molecular Biology and Evolution}
  \bibinfo{volume}{39}, \bibinfo{pages}{msac062}.
\newblock \DOIprefix\doi{10.1093/molbev/msac062}.
%Type = Article
\bibitem[{Zheng et~al.(2021)Zheng, Wang, Lynch and Gao}]{zheng2021}
\bibinfo{author}{Zheng, W.}, \bibinfo{author}{Wang, C.},
  \bibinfo{author}{Lynch, M.}, \bibinfo{author}{Gao, S.}, \bibinfo{year}{2021}.
\newblock \bibinfo{title}{The compact macronuclear genome of the ciliate
  {Halteria grandinella}: a transcriptome-like genome with 23,000
  nanochromosomes}.
\newblock \bibinfo{journal}{mBio} \bibinfo{volume}{12},
  \bibinfo{pages}{e01964--20}.
\newblock \DOIprefix\doi{10.1128/mBio.01964-20}.

\end{thebibliography}


\end{document}
